\providecommand{\zerJ}{\mathcal{J}}
\providecommand{\zerK}{\mathcal{K}}
\providecommand{\zersinc}{\operatorname{sinc}}
\providecommand{\zergE}{\gamma_{\mathrm{E}}}
\providecommand{\zerLL}{\mathcal{L}}
\providecommand{\zerconv}{\operatorname{conv}}
\providecommand{\zerdiag}{\operatorname{diag}}
\providecommand{\zerSym}{\mathfrak{D}}
\providecommand{\zerLog}{\operatorname{Log}}
\providecommand{\anapos}[1]{\bigl(#1\bigr)_{+}}
\providecommand{\ananeg}[1]{\bigl(#1\bigr)_{-}}
\providecommand{\anaone}{\mathbf{1}}
\providecommand{\anaE}{\mathcal{E}}
\providecommand{\anaEx}{\mathbb{E}}
\providecommand{\anaPr}{\mathbb{P}}
\providecommand{\anaM}{\mathcal{M}}
\providecommand{\anaW}{\mathcal{W}}
\providecommand{\anaZ}{\mathcal{Z}}
\providecommand{\anaC}{\mathcal{C}}
\providecommand{\filStir}[2]{\genfrac\{\}{0pt}{}{#1}{#2}}
\providecommand{\zfgdL}{d_{\mathrm L}}
\providecommand{\zfgE}{\mathcal E}
\providecommand{\zfgconv}{\operatorname{conv}}
\providecommand{\zfgdiam}{\operatorname{diam}}
\section{Zero-free intervals and positivity for trigonometric and algebraic polynomials}\label{sec:zfg}

This section treats extremal problems for positive measures and for real trigonometric and algebraic
polynomials with prescribed zeros or prescribed zero-free intervals. For a finite set of integer
frequencies, the longest interval on which a real trigonometric polynomial with that spectrum keeps one
sign is, after scaling, the least support diameter of a probability measure whose Fourier transform
vanishes on the spectrum (\cref{zfg:prop:duality}). On the side of measures this is the problem of the
shortest positive average that cancels prescribed zero contributions; the positive zero filters for zeros
of Dirichlet $L$-functions built on it are in \cref{sec:fil}, which uses the two-frequency classification
of \cref{fil:sec:two}. The questions treated here come from three sources.
\begin{itemize}
\item Kozma and Oravecz \cite{KozmaOravecz2003Gaps} determined the longest zero-free interval for a
progression of frequencies whose step is less than twice its first term. Their Question~1 asks for the
wide-step case, and their Question~2 asks whether the uniform bound $O(1/N)$ for the squares
$N^2,\dots,(N+K)^2$ has the right order. Two frequencies are classified in \cref{fil:sec:two}. Three-term
progressions reduce to a finite algebraic rule, with closed analytic branches for odd-integer and
half-odd-integer step ratios and for an infinite four-term family (\cref{zfg:sec:prog}). The order of
magnitude for squares is determined in \cref{zfg:sec:squares}.
\item Steinerberger \cite{Steinerberger2026Open} asked whether a nonnegative function on the circle
supported in $(-r,r)$, with vanishing Fourier coefficients on a finite set $A$, must have $r$ at least a
constant multiple of $(\sum_{a\in A}a^{-2})^{1/2}$. It must, with optimal constant $1/(\pi\sqrt2)$; this is a case
of sharp reciprocal-moment inequalities for the zero heights of a positive exponential sum
(\cref{zfg:sec:recip}).
\item Erd\'elyi \cite{Erdelyi1991Estimates} asked whether a polynomial of degree $n$ with all zeros in the
open ellipse with semiaxes $1$ and $\eps$ has Lorentz degree at least $cn/\eps^2$. It need not: the least
finite Lorentz degree has order $\max\set{n/\eps,\eps^{-2}}$ (\cref{zfg:sec:lorentz}).
\end{itemize}
All results are unconditional except the applications to zeros of Dirichlet $L$-functions marked \stG.

\subsection{Zero-free intervals and annihilating measures}\label{zfg:sec:dual}
\sources{ProveIt_New_Results_2026-11-27; positive-fourier-zeros}

For a finite nonempty set $A$ of positive integers let $M(A)$ be the supremum of the lengths of intervals of
$\R/\Z$ on which some nonzero real trigonometric polynomial
\[
 P(x)=\sum_{a\in A}\bigl(c_a\cos(2\pi ax)+d_a\sin(2\pi ax)\bigr)
\]
has no zero; coefficients may vanish, so the spectrum is only required to lie in $\pm A$. A zero-free real
function has constant sign, so $M(A)$ is also the supremum of the lengths of intervals on which such a $P$ is
strictly positive. For a finite set $A$ of positive reals let $B(A)$ be the least diameter of the support of a
probability measure $\mu$ on $\R$ with $\widehat\mu(a)=\int e^{iat}\dd\mu(t)=0$ for every $a\in A$.

\begin{proposition}[Duality; \stU]\label{zfg:prop:duality}
The minimum defining $B(A)$ exists, and for a finite set $A$ of positive integers $M(A)=B(A)/(2\pi)$. If
$A=\set{a_1<\dots<a_n}$, then
\[
 \max_{1\le m\le n}\frac{m^2}{2\sum_{j=1}^ma_j}\ \le\ M(A)\ \le\ \frac12\sum_{j=1}^n\frac1{a_j}.
\]
For every positive integer $g$, $M(gA)=M(A)/g$.
\end{proposition}

\begin{proof}
Convolving the measures $\tfrac12(\delta_0+\delta_{\pi/a})$, $a\in A$, gives an annihilating probability measure of
diameter $\pi\sum_{a\in A}a^{-1}$. Translating a minimizing sequence to start at $0$ puts it in a common compact
interval; a weakly convergent subsequence keeps the finitely many Fourier conditions, and its limit has
support of diameter at most the infimum. For a compact interval $I$ put $v_A(t)=(\cos(at),\sin(at))_{a\in A}$. An
annihilating probability measure on $I$ exists exactly when $0$ lies in the convex hull of $v_A(I)$. If
$\abs I<B(A)$, strict separation gives a linear functional positive on $v_A(I)$, that is, a polynomial with
spectrum in $\pm A$ that is strictly positive on $I$. Conversely, every interval longer than $B(A)$ contains a
translate of the support of a minimizer, against which every such polynomial integrates to zero, so none keeps
a strict sign there. For integer frequencies $B(A)\le2\pi$ (uniform measure on $[0,2\pi]$), so intervals on the
circle and on the line give the same supremum, and $t=2\pi x$ gives $M(A)=B(A)/(2\pi)$.

An annihilator of $A$ annihilates every prefix $\set{a_1,\dots,a_m}$, and \cref{fil:thm:meanbound} gives
$B(\set{a_1,\dots,a_m})\ge\pi m^2/\sum_{j\le m}a_j$; this is the lower bound. The convolution above gives the upper
bound, which is the one-dimensional case of the bound of Kozma and Oravecz \cite{KozmaOravecz2003Gaps}. Finally a
polynomial with spectrum in $\pm gA$ is $Q(gx)$ with $Q$ of spectrum in $\pm A$; it has no zero on an interval of
length $\ell$ exactly when $Q$ has none on an interval of length $g\ell$, and $Q$, having mean zero, vanishes on
every interval of length $1$.
\end{proof}

For integer spectra it is convenient to pass to the cosine coordinate. Let $S$ be a finite nonempty set of
positive integers, $T_n(\cos\phi)=\cos(n\phi)$, $V_S(x)=(T_n(x))_{n\in S}$, and $C_a=\zfgconv V_S([a,1])$ for
$-1\le a\le1$.

\begin{lemma}[Contact form; \stU]\label{zfg:lem:contact}
Let $a_*=\max\set{a:0\in C_a}$. Then $M(S)=\arccos(a_*)/\pi$, and there is $h\ne0$ with $h\cdot V_S\ge0$ on
$[a_*,1]$. The origin is a convex combination of at most $\abs S$ distinct contact vectors $V_S(x)$ with
$h\cdot V_S(x)=0$, and every such combination with positive weights uses $x=a_*$.
\end{lemma}

\begin{proof}
After the centre of an arc is translated to $0$, replacing $f(t)$ by $\tfrac12(f(t)+f(-t))$ preserves strict
positivity, and $x=\cos(2\pi t)$ maps an arc of length $\ell$ centred at $0$ onto $[\cos(\pi\ell),1]$. So a polynomial
strictly positive on such an arc amounts to a vector $h$ with $h\cdot V_S>0$ on $[\cos(\pi\ell),1]$, that is, to a
strict separation of $0$ from $C_{\cos(\pi\ell)}$. The uniform measure on the circle shows $0\in C_{-1}$; the sets
$C_a$ are compact and vary continuously in the Hausdorff distance, and $V_S(1)=(1,\dots,1)$, so the maximum exists
and gives the gap. If $0$ were interior to $C_{a_*}$, every direction would have support function bounded below by
a common positive number, and this would persist for slightly larger $a$; so a supporting hyperplane
$h^\perp$ passes through $0$. A representing measure lives on its contacts, and Carath\'eodory's theorem inside
the hyperplane bounds the number of nodes by $\abs S$. A positive representation using only nodes larger than
$a_*$ would contradict maximality.
\end{proof}

\begin{corollary}[Rational or transcendental; \stU]\label{zfg:cor:dichotomy}
For every finite set $S$ of positive integers, $M(S)$ is rational or transcendental.
\end{corollary}

\begin{proof}
The condition $0\in C_a$ is a first-order formula over $\Q$ in $a$, with at most $\abs S+1$ nodes in $[a,1]$,
nonnegative weights of sum one and the moment equations. By quantifier elimination for real closed fields
$a_*$ is algebraic, so $z=e^{i\pi M(S)}$, a root of $z^2-2a_*z+1$, is algebraic. If $M(S)$ were algebraic and
irrational, the Gelfond--Schneider theorem would make $z=(-1)^{M(S)}$ transcendental.
\end{proof}

\subsection{Two frequencies in arbitrary ratio}\label{fil:sec:two}
\sources{grh_positive_spectral_extensions_2026-11-27; dirichlet_zero_filters; grh_extensions_2026-11-27_2}

\Cref{rec:thm:hblock} prescribes consecutive harmonics. For two frequencies in arbitrary ratio, possibly
irrational, the problem is solved completely. Here $\widehat\mu(\omega)=\int e^{i\omega u}\dd\mu(u)$ for a compactly
supported probability measure $\mu$, and for $r>1$
\[
 \mathcal B(r)=\inf\set{\operatorname{diam}\supp\mu:\ \widehat\mu(1)=\widehat\mu(r)=0}.
\]
No symmetry of $\mu$ is assumed. For frequencies $0<\alpha<\gamma$ the least diameter is $\alpha^{-1}\mathcal B(\gamma/\alpha)$,
attained by the rescaled minimizers below.

\begin{theorem}[Exact two-frequency classification; \stU]\label{rec:thm:twofreq}
Let $r>1$.
\begin{enumerate}
\item If $r=n\ge3$ is an odd integer, then $\mathcal B(n)=\pi$, and $\tfrac12(\delta_{-\pi/2}+\delta_{\pi/2})$ is the unique
minimizer up to translation.
\item Otherwise let $n=2k+1<r<n+2$ with $k\in\Z_{\ge0}$ and $s=(-1)^k$. If $k=0$ put $t=0$ and $M=1$; if $k\ge1$ let $t$
be the unique solution in $(k\pi/r,(k+\tfrac12)\pi/r)$ of $\tan t=r\tan(rt-k\pi)$ and $M=s\cos(rt)/\cos t>0$.
There is a unique $T\in(\pi/2,(n+2)\pi/(2r))$ with $s\cos(rT)/\cos T=M$. Then $\mathcal B(r)=2T$, and with
\[
 P=\frac{-\cos T}{\cos t-\cos T},\qquad Q=\frac{\cos t}{\cos t-\cos T},
\]
the unique minimizer up to translation is $\mu_r=\tfrac P2(\delta_{-t}+\delta_t)+\tfrac Q2(\delta_{-T}+\delta_T)$, the inner
atoms combined when $t=0$.
\item For $1<r<3$ this is \cref{fil:cor:pair} with $\alpha=1$, $\gamma=r$, $m=1$: $T=2\pi/(1+r)$,
$\mathcal B(r)=4\pi/(1+r)$, $P=-c/(1-c)$, $Q=1/(1-c)$, $c=\cos(2\pi/(1+r))$.
\end{enumerate}
\end{theorem}

\begin{proof}
Let $r$ be nonresonant and $f(u)=s\cos(ru)/\cos u$ on $[0,\pi/2)$; since $s\cos(r\pi/2)<0$, $f\to-\infty$ at
$\pi/2$, so $\max f$ is attained. For $k=0$, $\cos(ru)\le\cos u$ on the only positive lobe $[0,\pi/(2r))$, so
$\max f=f(0)=1$, attained only at $0$. For $k\ge1$ a point of an earlier positive lobe moves by a multiple
of $2\pi/r$ into the last lobe $J=((k-\tfrac12)\pi/r,(k+\tfrac12)\pi/r)$ with the same numerator and a smaller
denominator, so the maximum lies in $J$. With $V(u)=s\{-r\sin(ru)\cos u+\cos(ru)\sin u\}$, $f'=V/\cos^2u$ and
$V'(u)=(1-r^2)s\cos(ru)\cos u<0$ on $J$; $V(k\pi/r)>0$ and $V<0$ at the right end of $J$, so there is exactly one
stationary point $t$, the unique maximizer, and $V(t)=0$ is the tangent equation. On
$(\pi/2,(n+2)\pi/(2r))$, $f>0$ runs from $+\infty$ to $0$ with $V'<0$ and $V(\pi/2^+)=s\cos(r\pi/2)<0$, so $f$ decreases
and $T$ exists uniquely.

Put $g(u)=M\cos u-s\cos(ru)$, even. On $[0,\pi/2)$, $g\ge0$ with equality only at $t$; on $(\pi/2,T)$,
$g=\cos u\,(M-f(u))>0$; $g(\pi/2)>0$. So $g\ge0$ on $[-T,T]$ with zero set $\set{\pm t,\pm T}$, double at $\pm t$
($g''(t)=(r^2-1)M\cos t>0$) and simple at $\pm T$. An admissible $\mu$ of diameter at most $2T$, translated into
$[-T,T]$, has $\int g\dd\mu=0$, so it lives on these nodes. The measure $\mu_r$ is admissible:
$P\cos t+Q\cos T=0$ and $\cos(ru)=sM\cos u$ at both nodes, and symmetry kills the sine moments. If
$t>0$ and $d_t,d_T$ are the differences of the masses at the positive and negative nodes, the sine
moments give $d_t\sin t+d_T\sin T=0=d_t\sin(rt)+d_T\sin(rT)$, and
$\sin t\sin(rT)-\sin T\sin(rt)=\sin t\{\sin(rT)-\tfrac{sM}r\sin T\}\ne0$, using $r\sin(rt)=sM\sin t$ and
$M\sin T-sr\sin(rT)>0$, which follows from $V(T)<0$ and $\cos T<0$. For $t=0$ the first sine moment
already balances the endpoints. Total mass and the first cosine moment then fix $P,Q$. Since $Q>0$
both endpoints carry mass, so no admissible measure is shorter. At an odd resonance $\cos u\ge0$ on
$[-\pi/2,\pi/2]$ forces the endpoints, the sine moment equal masses, and this measure annihilates the
odd frequency $n$; any shorter interval has positive cosine after centring. For $1<r<3$,
$\cos u-\cos(ru)=2\sin\frac{(1+r)u}2\sin\frac{(r-1)u}2$ gives (iii), which is also \cref{fil:cor:pair}.
\end{proof}

For $r=(p+1)/p$ this is the case $H=2$, one real part, of \cref{rec:thm:hblock}. The first-band formula and numerical two-mode solutions are known in the
input-shaping literature
\cite{CrainSinghoseSeering1996Derivation,SinghoseCrainSeering1997Convolved,Convolve1995Minimizing}, and the
setting belongs to the extremal theory of positive-definite bandlimited functions \cite{Logan1983Extremal}.
The new feature is the behaviour at the odd resonances, which is asymmetric. Already the coalescence of the
two frequencies is singular.

\begin{corollary}[Merging frequencies; \stU]\label{fil:cor:merger}
A single frequency imposed once needs diameter $\pi$, but $\lim_{r\downarrow1}\mathcal B(r)=2\pi$. The minimizers
for $r\downarrow1$ converge weakly to $\tfrac14\delta_{-\pi}+\tfrac12\delta_0+\tfrac14\delta_\pi$, whose transform
$\cos^2(\pi\omega/2)$ has a double zero at $\omega=1$; it is the unique shortest measure with
$\widehat\mu(1)=\widehat\mu'(1)=0$, by \cref{rec:thm:fwidth}.
\end{corollary}

\begin{proof}
Let $r\downarrow1$ in \cref{fil:cor:pair} with $\alpha=1$, $\gamma=r$, $m=1$, that is, in
\cref{rec:thm:twofreq}(iii): $T\to\pi$, $c\to-1$, $P\to\tfrac12$, $Q\to\tfrac12$.
\end{proof}

\begin{theorem}[Asymmetric resonance law; \stU]\label{rec:thm:cusp}
Fix an odd $n\ge3$ and put $A_n=n(n^2-1)/6$, $C_n=(3\pi/(2n(n^2-1)))^{1/3}$. As $\eps\downarrow0$, for $r=n+\eps$,
\[
 t=\tfrac\pi2-C_n\eps^{1/3}+O_n(\eps),\qquad T=\tfrac\pi2+2C_n\eps^{1/3}+O_n(\eps),\qquad
 \mathcal B(n+\eps)=\pi+4C_n\eps^{1/3}+O_n(\eps),
\]
with $P=\tfrac23+O_n(\eps^{2/3})$ and $Q=\tfrac13+O_n(\eps^{2/3})$. With $s_n=(-1)^{(n-1)/2}$ and the finite
maximum $M_n^-=\max_{0\le u<\pi/2}(-s_n\cos(nu)/\cos u)$,
\[
 \mathcal B(n-\eps)=\pi+\frac{\pi}{n+M_n^-}\eps+O_n(\eps^2).
\]
In particular $M_3^-=1$, $\mathcal B(3-\eps)=\pi+\pi\eps/4+O(\eps^2)$, and $M_5^-=-\min_{0\le v\le1}(16v^4-20v^2+5)=\tfrac54$,
$\mathcal B(5-\eps)=\pi+\tfrac{4\pi}{25}\eps+O(\eps^2)$.
\end{theorem}

\begin{proof}
Put $h=\pi\eps/2$ and $t=\pi/2-y$, so $M=\sin(ry-h)/\sin y$. At $\eps=0$ the quotient $\sin(ny)/\sin y$, a sum of
$n$ unit vectors with distinct adjacent phases, is below $n$ for $0<y\le\pi/2$, while the value at
$y=\eps^{1/3}$ tends to $n$; hence $y\to0$. The stationary equation
$H(\eps,y)=r\cos(ry-h)\sin y-\sin(ry-h)\cos y=0$ expands as
$\tfrac\pi2\eps-2A_ny^3+O_n(y^5+\eps y^2+\eps^2y+\eps^3)$. With $z=\eps^{1/3}$ and $y=zv$, $H(z^3,zv)/z^3$ is
real analytic at $z=0$ and equals $\tfrac\pi2-2A_nv^3+O_n(z^2)$, with $v$-derivative $-6A_nC_n^2\ne0$ at
$(0,C_n)$, so $v=C_n+O_n(z^2)$. For the outer node $T=\pi/2+x$ the crossing equation reads
$\sin(rx+h)/\sin x=\sin(ry-h)/\sin y$, and $x\to0$ as before. With $x=zw$,
\[
 \frac{\sin((n+z^3)zw\pm\pi z^3/2)}{\sin(zw)}=n+z^3+z^2\Bigl(\pm\frac\pi{2w}-A_nw^2\Bigr)+O_n(z^4),
\]
so after subtracting the two sides and dividing by $z^2$ the equation becomes
$A_n(2C_n^3/w-w^2+3C_n^2)+O_n(z^2)=0$; since $w^3-3C_n^2w-2C_n^3=(w-2C_n)(w+C_n)^2$ and the derivative at
$2C_n$ is $-\tfrac92A_nC_n\ne0$, $w=2C_n+O_n(z^2)$, and uniqueness of $T$ identifies this root. The weights
follow from $\cos t=\sin y$ and $-\cos T=\sin x$. For $r=n-\eps$ the sign of the band is $-s_n$, the limiting
maximum $M_n^-$ is attained before $\pi/2$ (at $0$ for $n=3$, at a nondegenerate interior point for $n\ge5$),
the maximum is $M_n^-+O_n(\eps)$, and with $x=\eps v$ the outer equation
$-\sin((n-\eps)x-h)/\sin x=M_n^-+O_n(\eps)$ becomes $\pi/(2v)-n=M_n^-$ at $\eps=0$, with nonzero derivative.
\end{proof}

A perturbation $\eps$ of the higher frequency just above an odd multiple of the lower one therefore
costs a window increase of order $\eps^{1/3}$, and just below it only of order $\eps$. For
$\eps(\nu)=\max\set{\abs{\widehat\nu(1)},\abs{\widehat\nu(r)}}$ and nonresonant $r$, with $t,T,M,P,Q,g$ as above:

\begin{proposition}[Sharp stability for two frequencies; \stU]\label{rec:prop:twostable}
Every probability measure $\nu$ on $[-T,T]$ has $W_1(\nu,\mu_r)\le C_r\sqrt{\eps(\nu)}$, and the exponent $\tfrac12$
cannot be increased. At an odd resonance, on $[-\pi/2,\pi/2]$, the optimal exponent is $1$.
\end{proposition}

\begin{proof}
With $G$ the zero set of $g$, $g(u)\ge c_r\operatorname{dist}(u,G)^2$ on $[-T,T]$ and $\int g\dd\nu\le(M+1)\eps(\nu)$, so
rounding to $G$ costs $O_r(\sqrt\eps)$ by Cauchy--Schwarz and moves the moments by as much; total mass and
the two complex moments determine weights on $G$ through the invertible linear algebra of the
uniqueness proof. Splitting an inner atom by $\pm a$ changes the moments by $O_r(a^2)$ at transport
cost $\asymp a$. At a resonance $\cos u$ grows linearly from the endpoints, and moving them inward proves
sharpness.
\end{proof}

For $1<r<3$ the first assertion is a case of \cref{fil:thm:stable}, at the strictly positive pair base
of \cref{fil:cor:pair} with $\gamma<3\alpha$.

\begin{theorem}[First-order cost of shortening at two frequencies; \stU]\label{rec:thm:twoshort}
For nonresonant $r$ and small $d>0$ let $D_r(d)=\inf\set{\eps(\nu):\supp\nu\subset[-T+d/2,T-d/2]}$ over
probability measures. Then $D_r(d)=L_rd+O_r(d^2)$ with $L_r=Q(-g'(T))/(2(M+1))>0$, and for odd $n\ge3$,
$L_{n+\eps}=\frac{3A_nC_n^2}{2(n+1)}\eps^{2/3}+O_n(\eps)$. At an odd resonance $r=n\ge3$, with $T=\pi/2$,
\[
 D_n(d)=\frac{n+M_n^-}{2(1+M_n^-)}\,d+O_n(d^2);
\]
thus $D_3(d)=d+O(d^2)$ and $D_5(d)=\tfrac{25}{18}d+O(d^2)$.
\end{theorem}

\begin{proof}
Put $a=-g'(T)/2$, $K=Qa$ and $h_0(u)=K\cos u/\cos t$, equal to $K$ at the inner nodes and $K-a$ at the outer
ones. On the shortened interval $g+dh_0\ge dK-O_r(d^2)$: near an inner node the quadratic minimum loses
only $O(d^2)$, at each moved endpoint the increase is $da$, and elsewhere $g$ is bounded below.
Integrating gives $(M+1+dK/\cos t)\eps(\nu)\ge dK-O_r(d^2)$. Conversely, keep the inner nodes, move the outer
ones to $\pm(T-d/2)$ and choose symmetric weights with $\widehat\nu(1)+s\widehat\nu(r)=0$, possible for small $d$ since
$\cos u+s\cos(ru)$ has opposite signs at the two nodes when $d=0$; then $(M+1)\widehat\nu(1)=\int g\dd\nu=dQa+O_r(d^2)$.
At $r=n+\eps$, $g(\pi/2+x)=h+3A_ny^2x-A_nx^3+O_n(\eps^{4/3})$ with $y=C_n\eps^{1/3}+O_n(\eps)$ and $x=2y+O_n(\eps)$,
which gives $-g'(T)=9A_nC_n^2\eps^{2/3}+O_n(\eps)$; use $Q=\tfrac13+O_n(\eps^{2/3})$ and $M=n+O_n(\eps^{2/3})$.

At $r=n$ put $\sigma=\sin(n\pi/2)$ and let $a_*<\pi/2$ be the unique maximizer of $-\sigma\cos(nu)/\cos u$ on
$[0,\pi/2)$ (the lobe argument above), with value $M=M_n^-$. Then $q(u)=M\cos u+\sigma\cos(nu)\ge0$ on
$[-\pi/2,\pi/2]$, with double zeros at $\pm a_*$ and $q'(\pi/2)=-(M+n)$. The perturbation
$q+\tfrac d2\,\frac{M+n}{\cos a_*}\cos u$ is at least $\tfrac d2(M+n)-O_n(d^2)$ on the shortened interval,
which gives $D_n(d)\ge\frac{M+n}{2(M+1)}d-O_n(d^2)$. For the upper bound put total mass
$\alpha_d=\frac{(n-1)d}{2(M+1)\cos a_*}$ on the pair $\pm a_*$ and $1-\alpha_d$ on $\pm(\pi/2-d/2)$; the two
real moments are $\frac{M+n}{2(M+1)}d+O_n(d^2)$ and $\sigma$ times this, using
$\cos(na_*)=-\sigma M\cos a_*$, and the imaginary moments vanish.
\end{proof}

\begin{corollary}[A second singularity at resonance; \stU]\label{fil:cor:kappa}
For odd $n\ge3$, $L_{n-\eps}\to L_n=\frac{n+M_n^-}{2(1+M_n^-)}$, while $L_{n+\eps}\sim\frac{n(n-1)}4C_n^2\eps^{2/3}\to0$.
For fixed $0<d<\pi$, $D_r(d)\to D_n(d)$ as $r\to n$, so the limits $d\downarrow0$ and $r\downarrow n$ of
$D_r(d)/d$ do not commute.
\end{corollary}

\begin{proof}
The left limit follows from the continuity of the maximizer used in \cref{rec:thm:cusp}; the right one is
\cref{rec:thm:twoshort}, since $3A_n/(2(n+1))=n(n-1)/4$. For fixed $d$, compactness of probability measures on
a common interval and linear rescaling of the support intervals give $D_r(d)\to D_n(d)$.
\end{proof}

\begin{example}\label{rec:ex:ratio4}
For $r=4$ put $x=\sqrt{(2+\sqrt{10})/12}$ and $z=-x+\sqrt{1-2x^2}$. Tangency of $-\cos(4u)/\cos u$ gives
$24x^4-8x^2-1=0$, and with $M=-(8x^4-8x^2+1)/x$ the identity $8v^4-8v^2+1+Mv=8(v-x)^2(v^2+2xv+3x^2-1)$ identifies
$z\in(-1,0)$ as the outer crossing, the other quadratic root being below $-1$. Hence $t=\arccos x$,
$T=\arccos z$ and
\[
 \mathcal B(4)=2\arccos\Bigl(-\sqrt{\tfrac{2+\sqrt{10}}{12}}+\sqrt{\tfrac{4-\sqrt{10}}6}\Bigr)\approx3.7138292065,
\]
with $P\approx0.3008471646$ and $Q\approx0.6991528354$.
\end{example}

The width is unchanged by a common real part: for $\beta\in\R$ and $0<\alpha<\gamma$ the least support diameter
of a probability measure with $\int e^{(\beta+i\alpha)u}\dd\mu=\int e^{(\beta+i\gamma)u}\dd\mu=0$ is $\alpha^{-1}\mathcal B(\gamma/\alpha)$, since the
tilt $e^{\beta u}\dd\mu/\int e^{\beta v}\dd\mu$ preserves the support; the optimal nodes $u_j$ and masses $w_j$ become
$\widetilde w_j=w_je^{-\beta u_j}/\sum_\ell w_\ell e^{-\beta u_\ell}$. Under \cref{src:thm:dgrh} two ordinates $\alpha<\gamma$ of the same or
of different Dirichlet $L$-functions lie on $\beta=\tfrac12$, and averaging $S_W$, or any explicit formula in which a
zero contributes $c_\rho x^\rho$, at the points $xe^{u_j}$ with the weights $\widetilde w_j$ cancels both contributions
exactly, with the least possible multiplicative scale ratio $\exp(\mathcal B(\gamma/\alpha)/\alpha)$ (\stG). Equal
vanishing orders at both frequencies are settled for $\gamma\le3\alpha$ by \cref{fil:cor:pair}. For the first two zeta ordinates
$\gamma_2/\gamma_1\approx1.48726$ lies in the first band, so the least logarithmic width is
$4\pi/(\gamma_1+\gamma_2)\approx0.357438$, with three positive weights $\approx0.30032$, $0.44851$, $0.25117$ at the
shifts $0$, $0.178719$, $0.357438$ (rounded):
\begin{center}
\artifact{grh_positive_spectral_extensions_2026-11-27/results/filter_checks.json}\\
\artifact{grh_positive_spectral_extensions_2026-11-27/results/cusp.csv}\\
\artifact{grh_extensions_2026-11-27_2/evidence/two_frequency_data.json}
\end{center}
The last file checks the classification at 21 ratios and the right cusp at $n=3,5,7,9$ down to
$\eps=10^{-18}$, at 90 digits; it is numerical, not an interval certificate.

\paragraph{Zero-free intervals and the position inside the convolution bound.}
The two-frequency problem is dual to the longest interval on which a trigonometric polynomial with
prescribed spectrum keeps one sign; this answers the two-frequency case of a question of Kozma and Oravecz
\cite{KozmaOravecz2003Gaps}.

\begin{corollary}[Gaps for two integer frequencies; \stU]\label{rec:cor:kozma}
For integers $0<N<M$ the supremum of the lengths of zero-free intervals on $\R/\Z$ of real trigonometric
polynomials with spectrum in $\set{\pm N,\pm M}$ is $\mathcal B(M/N)/(2\pi N)$; for $M=3N$ it is $1/(2N)$.
\end{corollary}

\begin{proof}
Scaling $u=Nt$ gives $B(\set{N,M})=\mathcal B(M/N)/N$; apply \cref{zfg:prop:duality}.
\end{proof}

In the variables $u=\pi(r-n)/2$, $t=\pi/2-(u+y)/r$, $T=\pi/2+z/r$, the equations of \cref{rec:thm:twofreq}(ii) read
$\tan y=r\tan((u+y)/r)$ and $\sin(u+z)=M\sin(z/r)$ with $M=\sin y/\sin((u+y)/r)$, $y\in(0,\pi/2)$, $z\in(0,\pi-u)$. Convolving the
two single-frequency minimizers gives $\mathcal B(r)\le\pi+\pi/r$.

\begin{theorem}[Ratios attaining the convolution bound; \stU]\label{rec:thm:twosat}
For each odd $n\ge3$ exactly one $r\in(n,n+2)$ has $\mathcal B(r)=\pi+\pi/r$. It lies in $(n,n+1)$, satisfies
$\tan(\pi(r-n)/2)=r^{-1}\cot(\pi/(2r))$, and $r-n\to(2/\pi)\arctan(2/\pi)$ as $n\to\infty$. Uniformly for $u$ in compact subsets of $(0,\pi)$,
\[
 \mathcal B(r)=\pi+2Z(u)/r+O(r^{-3}),
\]
where $Y(u)\in(0,\pi/2)$ solves $\tan Y=u+Y$ and $Z(u)$ is the first positive root of $\sin(u+Z)=Z\cos Y$.
\end{theorem}

\begin{proof}
The convolution has nodes $\pm(\pi/2\pm\pi/(2r))$, so equality forces the inner node to be the unique contact,
$u+y=\pi/2$, and the contact equation becomes the criterion; conversely the criterion makes $y=\pi/2-u$ the
contact and $z=\pi/2$ a root of the crossing equation. With $v=\pi/(2r)$, $\frac{\dd}{\dd r}(r^{-1}\cot v)=(v-\sin v\cos v)/(r^2\sin^2v)\le\pi^2v/(6r^2)<\pi/2$,
using $v-\sin v\cos v\le2v^3/3$ and $\sin v\ge2v/\pi$, while $\tan(\pi(r-n)/2)$ increases from $0$ to $\infty$ on $(n,n+1)$ with
derivative at least $\pi/2$, and the phase lies outside $(0,\pi/2)$ on $(n+1,n+2)$; and $r^{-1}\cot(\pi/(2r))\to2/\pi$. For the
expansion, $r\tan((u+y)/r)=u+y+O(r^{-2})$ and the limiting equation has derivative $\sec^2Y-1>0$, so $y=Y+O(r^{-2})$;
the crossing equation, written $\sin(u+z)=z(M/r)\sin(z/r)/(z/r)$, is a smooth perturbation of its limit, whose
derivative at the root is negative since $\sin(u+z)/z$ decreases before $z=\pi-u$.
\end{proof}

Under \cref{src:thm:dgrh}, for positive ordinates $\alpha<\gamma$ the least logarithmic span of a positive average
cancelling $\tfrac12+i\alpha$ and $\tfrac12+i\gamma$ is $d(\alpha,\gamma)=\alpha^{-1}\mathcal B(\gamma/\alpha)\in[\pi/\alpha,\pi/\alpha+\pi/\gamma]$. Put
$E_\alpha(\gamma)=\gamma(d(\alpha,\gamma)-\pi/\alpha)=r(\mathcal B(r)-\pi)\in[0,\pi]$, $r=\gamma/\alpha$.

\begin{theorem}[A universal law for the extra width; \stG]\label{rec:thm:twolaw}
Fix an ordinate $\alpha>0$ of some Dirichlet $L$-function and a primitive $\chi$, possibly of conductor one, and let
$N_{\chi,\alpha}(T)$ count the ordinates $\alpha<\gamma\le T$ of $L(s,\chi)$ with multiplicity. For every continuous $f$ on $[0,\pi]$,
\[
 \lim_{T\to\infty}\frac1{N_{\chi,\alpha}(T)}\sum_{\alpha<\gamma\le T}f(E_\alpha(\gamma))=\frac1\pi\int_0^\pi f(2Z(u))\dd u.
\]
The law $X$ of $2Z(U)$, $U$ uniform on $(0,\pi)$, has no atoms and support $[0,\pi]$, so $\liminf E_\alpha(\gamma)=0$ and
$\limsup E_\alpha(\gamma)=\pi$. With $Y_\pi\in(0,\pi/2)$ solving $\tan Y_\pi=\pi+Y_\pi$ and $t_0=\pi/2$, as $e\downarrow0$,
\[
 \Pr(X\le e)\sim\frac{1+\cos Y_\pi}{2\pi}\,e,\qquad\Pr(X\ge\pi-e)\sim\frac2\pi\frac{\sqrt{2t_0^3}}{1+t_0^2}\sqrt e.
\]
\end{theorem}

\begin{proof}
The positive ordinates of $L(s,\chi)$ are equidistributed modulo one after multiplication by any fixed nonzero
real \cite[Cor.~8]{AkbaryMurty2008Uniform}, so the phases $u_\gamma=\pi\set{\gamma/(2\alpha)-\tfrac12}$ are equidistributed on $[0,\pi)$; this is
unconditional, and \cref{src:thm:dgrh} places all multipliers on one vertical line. On $[\eta,\pi-\eta]$,
\cref{rec:thm:twosat} gives $E_\alpha(\gamma)=2Z(u_\gamma)+O_\eta(r^{-2})$; uniform continuity and $\eta\downarrow0$ give the limit. As $u\uparrow\pi$, $Z\to0$,
and at $u_*=\arctan(2/\pi)$ one has $Y=\pi/2-u_*$, $Z=\pi/2$, so the support is $[0,\pi]$; $Z$ is real analytic and
nonconstant, so level sets are null. At $u\downarrow0$, $u=\tan Y-Y\sim Y^3/3$ and $Z\sim2(3u)^{1/3}$; at $u=\pi-v$, $Z\sim v/(1+\cos Y_\pi)$,
which gives the lower edge. A stationary point of $Z$ satisfies $\cot(u+Z)=-\cot Y$, hence $Z=\tan Y=u+Y$, forcing
$Z=t_0$ and $u=u_*$; there $Z''(u_*)=-(1+t_0^2)^2/(2t_0^3)$, and the set where $2Z\ge\pi-e$ has length $\sim2\sqrt{e/\abs{Z''(u_*)}}$.
\end{proof}

Numerically $\Pr(X\le\pi/2)=0.328039432\ldots$, and the median of $X/\pi$ is $0.711649171\ldots$
(\artifact{dirichlet_zero_filters/law_results.txt}).

\begin{theorem}[Arithmetic second term; \stG]\label{rec:thm:twobias}
Fix $\alpha>0$, a primitive $\chi$ and a smooth $\varphi$ supported in $(0,\pi)$; put $\kappa=1/(2\alpha)$, $H_\varphi(t)=\varphi(2Z(\pi\set{t-1/2}))$, and let
$N_\chi^+(T)$ count the zeros with $0<\gamma\le T$. Then
\[
 \sum_{\alpha<\gamma\le T}\varphi(E_\alpha(\gamma))=N_\chi^+(T)\int_0^1H_\varphi+T\int_0^1H_\varphi(t)g_{\chi,\kappa}(t)\dd t+o(T),
\]
where $g_{\chi,\kappa}=0$ unless $\kappa=A\log p/(2\pi B)$ for a prime $p$ and coprime positive integers $A,B$, and then
$g_{\chi,\kappa}(t)=-\frac{\log p}\pi\Re\sum_{j\ge1}\chi(p)^{Aj}e^{-2\pi iBjt}p^{-Aj/2}$. A prime dividing the conductor gives no correction.
\end{theorem}

\begin{proof}
Away from odd ratios $E_\alpha(\gamma)=2z$ with band phase $u=\pi\set{\kappa\gamma-1/2}$, and $2z=2Z(u)+O(r^{-2})$ on compact phase intervals;
both profiles tend to $0$ uniformly at the seam (as $u\to\pi$, $0<z<\pi-u$; as $u\to0$, $y\to0$, $M/r\to1$ and $z\to0$), so
for large $r$ both tests vanish near the seam and differ by $O(r^{-2})$ elsewhere, and $\sum\gamma^{-2}<\infty$. Apply
\cite[Thm.~5, (3.2)]{FordSoundararajanZaharescu2009Distribution} to the smooth periodic $H_\varphi$; the orientation for complex
characters matches Landau's formula at $x=p^{Aj}$ \cite{Landau1911Nullstellen}.
\end{proof}

For $\alpha=\gamma_1$ and the zeta ordinates $\gamma_{10}$, $\gamma_{20}$ (ratios $3.521387$, $5.457824$) the joint widths are $0.283311228$ and
$0.262502497$, against $0.285377967$ and $0.262983914$ for the two separate filters
(\artifact{dirichlet_zero_filters/pair_results.txt}).

\paragraph{Wide arithmetic progressions.}
For positive integers $N,b$ and $K\ge0$ let $M_K=M_K(N,b)$ be the supremum of the lengths of intervals of $\R/\Z$ on which
a real trigonometric polynomial with spectrum in $\set{\pm(N+jb):0\le j\le K}$ is strictly positive. For $b<2N$,
$M_K=(K+1)/(2N+Kb)$ \cite[Thm.~3]{KozmaOravecz2003Gaps}, and the wide-step case was asked there.

\begin{theorem}[Wide-step limit and sharp order; \stU]\label{rec:thm:wide}
Let $b\ge2N$, $\beta=N/b$ and $c=\ceil{1/(2\beta)}$. If $b/(2N)$ is an integer then $bM_K=c$ for every $K$. Otherwise
$(2\beta c-1)/(K+2\beta)\le c-bM_K$ for every $K\ge0$, and $c-bM_K\le C_\beta/K$ for large $K$. In both cases $M_K\to\ceil{b/(2N)}/b$.
\end{theorem}

\begin{proof}
In $t=2\pi bx$ the allowed functions are $\Re(e^{i\beta t}P(e^{it}))$ with $\deg P\le K$. At $t=2\pi j$, $0\le j\le c$, every mode equals
$e^{ij\theta}$, $\theta=2\pi\beta$, and these points have $0$ in their convex hull (all circular gaps are at most $\pi$), so a
probability measure on them annihilates the whole progression and $bM_K\le c$; at resonance the lowest cosine
attains it. Otherwise $c\ge2$; on a positive interval of length $2\pi(c-\eps)$, $0\le\eps<1$, $P$ has no root on the circle,
the argument of $e^{i\beta t}P(e^{it})$ changes by less than $\pi$, an exterior root $w$ contributes at least $-\pi\eps$ since
$\frac{\dd}{\dd t}\arg(e^{it}-w)<\tfrac12$ and the change over $c$ full turns is zero, and an interior root contributes a
nonnegative amount; so $2\pi\beta(c-\eps)-K\pi\eps<\pi$, the lower bound. For the upper bound let $\alpha'=2\beta<1$,
$H_\rho(t)=e^{i\beta t}(1-\rho e^{it})^{-\alpha'}$ and $P_K(z)=\sum_{j\le K}\frac{(\alpha')_j}{j!}\rho^jz^j$. The argument $\psi_\rho$ of $H_\rho$ has
$\psi_\rho'=\beta(1-\rho^2)/\abs{1-\rho e^{it}}^2>0$ and $\psi_\rho(t+2\pi)=\psi_\rho(t)+2\pi\beta$, so after rotation by $e^{-i\pi c\beta}$ its phase on
$[\delta,2\pi c-\delta]$ stays within $\pi/2-m_0$ of $0$, with
$m_0=\pi(\tfrac12-(c-1)\beta)-\alpha'\arctan\bigl(\frac{1-\rho}{1+\rho}\cot\frac\delta2\bigr)$. The relative truncation error on $\abs z=\rho$ is at most
$2(K+2)^{\alpha'-1}\rho^{K+1}(1-\rho)^{\alpha'-1}$. Choosing $\rho=1-D/(K+2)$ and $\delta=U/(K+2)$ with $D$ and then $U/D$ large gives a
positive real part on the interval and $bM_K\ge c-U/(\pi(K+2))$.
\end{proof}

By the convex-hull duality of \cref{zfg:prop:duality}, the infinite progression has a shortest annihilating
probability measure of diameter exactly $c/b$ in circle units. Rational enclosures for nine instances, for
example $1.98724292\le bM_{256}\le1.99766356$ at $(N,b)=(2,5)$, are in \artifact{dirichlet_zero_filters/progression_results.txt}.
Exact values for $K=2$ and every $(N,b)$, and for $K=3$ when $b/N$ is an odd integer, are in \cref{zfg:sec:prog}.

\begin{theorem}[A growing collection of zeros; \stG]\label{rec:thm:twopack}
Fix $N,b\ge1$ with $b\ge2N$, a primitive $\chi$, $p_0>1$ and $0<\kappa<1/(2p_0)$, and let $K=K(T)\to\infty$ with
$K\le\kappa\log T/\log\log T$. There is $H_\chi$ such that for large $T$ every interval $[T(N+bj),T(N+bj)+H_\chi]$, $0\le j\le K$, contains
an ordinate $\gamma_j$ of $L(s,\chi)$, and for every such choice the least logarithmic span $D_T$ of a positive average
cancelling all $\tfrac12+i\gamma_j$ satisfies $\abs{TD_T-2\pi M_K}\le K^{-p_0}$. In particular $2\pi c/b-TD_T\asymp_{N,b}K^{-1}$ if $c>b/(2N)$; one
may take $p_0=2$ for $\kappa<1/4$.
\end{theorem}

\begin{proof}
The count $N^+(X)=X\log X/(2\pi)+c_\chi X+O_\chi(\log X)$ \cite{Selberg1992Old,FordSoundararajanZaharescu2009Distribution} gives the
$\gamma_j$. With $\lambda_j=N+bj$ and $\lambda'_j=\gamma_j/T$, tilting gives $TD_T=W(\lambda')$ and duality $W(\lambda)=2\pi M_K$, where $W$ is the least
width annihilating the frequencies. Let $d=K+1$, $m'=2d+1$, $\delta=K^{-p_0}$. An optimal measure for $\lambda$ has at most $m'$
atoms (Carath\'eodory), one of mass at least $1/m'$; convolving with the uniform law on $j\delta/(m'-1)$, $0\le j<m'$,
costs width $\delta$, keeps the annihilation, gives mass at least $m'^{-2}$ at each new point, and leaves target
moments $O(H_\chi/T)$. The Vandermonde matrix of the reserved weights has nodes at distances $\gg\Delta\abs{\ell-h}$,
$\Delta=\delta/(m'-1)$, so the entries of its inverse are at most $(C'/\delta)^{2d}$, and the real correction has size
$T^{-1+2p_0\kappa+o(1)}<m'^{-2}$; so $W(\lambda')\le W(\lambda)+\delta$, and symmetrically. The last claim follows from \cref{rec:thm:wide}.
\end{proof}




\subsection{Three and four frequencies in a wide progression}\label{zfg:sec:prog}
\sources{positive-fourier-zeros; kozma_oravecz_147_2026-11-27; kozma_oravecz_three_frequency_2026-11-28; kozma_oravecz_half_integer_2026-11-28; proveit_three_extensions_2026-11-27}

In the notation of \cref{rec:thm:wide}, $M_K(N,b)=M(\set{N+jb:0\le j\le K})$. If $b/(2N)$ is an integer then
$M_K(N,b)=b^{-1}\ceil{b/(2N)}=1/(2N)$ for every $K$ by \cref{rec:thm:wide}. This subsection determines $M_2(N,b)$ for all $N,b$ by
a finite algebraic rule, gives closed analytic branches when $b/N$ is an odd integer or half an odd integer at
least $5/2$, and does the same for $K=3$ when $b/N$ is an odd integer.

\paragraph{A finite rule for three frequencies.}
By \cref{zfg:prop:duality} one may assume $\gcd(N,b)=1$. Put $m=N+b$,
\[
 V(x)=(T_N(x),T_m(x),T_{N+2b}(x)),\qquad D(x)=V(x)\times V'(x),
\]
and call a set of nodes in $[-1,1]$ feasible if the origin lies in the convex hull of its images under $V$; for at
most three nodes this is linear algebra and sign comparisons.

\begin{theorem}[A finite rule for three frequencies; \stU]\label{zfg:thm:threerule}
Form the following finite list of feasible node sets.
\begin{enumerate}
\item Every feasible singleton or pair with nodes in the union of the grids $\set{\cos(j\pi/b):0\le j\le b}$ and
$\set{\cos((2j+1)\pi/(2m)):0\le j<m}$.
\item Let $E$ consist of $1$ and the common real zeros of the coordinates of $D$ at which $V\ne0$. For each nonzero
normal $h$, up to scaling, of one of the forms
\[
\begin{aligned}
 &V(u)\times V(v), && u,v\in E,\\
 &D(v), && u\in E,\quad D(v)\cdot V(u)=0,\quad V(u)\times V(v)\ne0,\\
 &D(u), && D(u)\times D(v)=0,\quad D(u),D(v)\ne0,\quad V(u)\times V(v)\ne0,
\end{aligned}
\]
every feasible triple among the distinct real roots of $h\cdot V(x)$ in $[-1,1]$.
\end{enumerate}
The largest smallest node among the listed sets is the number $a_*$ of \cref{zfg:lem:contact}, so
$M_2(N,b)=\arccos(a_*)/\pi$. The list requires univariate root computations and finitely many isolated solutions of
bivariate polynomial systems.
\end{theorem}

\begin{proof}
Every listed set is feasible, so its smallest node is at most $a_*$. Choose a positive representation of $0$ on
$[a_*,1]$ with the fewest nodes; by \cref{zfg:lem:contact} it has at most three and uses $a_*$.

A singleton is a common zero of the coordinates, hence a zero of $T_m$, in the second grid. For a pair write the
nodes as $\cos\phi,\cos\theta$ with $0\le\phi<\theta\le\pi$; their vectors are antiparallel. If $T_m$ vanishes at
$\cos\theta$ it vanishes at $\cos\phi$, and both nodes lie in the second grid. Otherwise the identity
$\cos(Nt)+\cos((N+2b)t)=2\cos(bt)\cos(mt)$ forces $\cos(b\theta)=\cos(b\phi)$. If $\sin(b\theta)=0$, both nodes lie in the
first grid. Otherwise write $b\phi\equiv\sigma b\theta\pmod{2\pi}$ with $\sigma=\pm1$. The first two coordinates and
$\sin(b\theta)\ne0$ determine the cosine and the sine of $N\phi$ up to the same negative factor as those of $\sigma N\theta$;
both being unit vectors, the factor is $-1$. Thus $d=\phi-\sigma\theta$ has $Nd/\pi$ odd and $bd/\pi$ even, so coprimality makes
$d/\pi$ an odd integer, with $N$ odd and $b$ even. Then $\theta\ge\pi/2$, while $x=0$ is a feasible singleton because all
three frequencies are odd; it is listed and at least as large.

For a minimal triple no two vectors are proportional: opposite vectors give a feasible pair, and positively
proportional ones can be merged after renormalizing the weights. The normal $h$ of the supporting plane through
$0$ is therefore determined by two of the vectors. Each of the two nodes larger than $a_*$ is the endpoint $1$, a
common zero of the coordinates of $D$, or an interior regular tangency, where $h$ is proportional to $D$. This
gives the three cases of the list.

It remains to see that the list is finite. The coordinates of $V$ have distinct degrees, so $h\cdot V\not\equiv0$ for
$h\ne0$. $D\not\equiv0$, since otherwise the projective ratios of the coordinates of $V$ would be constant. For fixed
$p=V(u)\ne0$, $D(v)\cdot p\not\equiv0$: in coordinates with $p=(0,0,1)$, identical vanishing says that the ratio of the
first two transformed coordinates has zero derivative, contrary to their linear independence. Suppose a curve of
pairs $(u,v)$ with $V(u)\times V(v)\ne0$ satisfied the conditions of the third case. On a smooth piece choose a
smooth normal $h$ with one coordinate equal to one. Differentiating $h\cdot V(u)=h\cdot V(v)=0$ along the piece and
using the tangencies gives $h'\cdot V(u)=h'\cdot V(v)=0$, so $h'$ is proportional to $h$ and the normalization gives
$h'=0$; but a fixed $h$ has finitely many contacts. The same argument applies on the smooth locus of a complex
component, so dividing the three minors of $D(u)\times D(v)$ by their polynomial greatest common divisor discards the
common curve components (for instance the diagonal $u=v$) without losing a valid pair, and the remaining system is
zero-dimensional.
\end{proof}

\paragraph{Contact certificates.}
The closed values below rest on one deduction.

\begin{lemma}[Contact certificate; \stU]\label{zfg:lem:cert}
Let $S$ be a set of $k+1$ positive integers, $0<\theta_1,\dots,\theta_k<\theta_0<\pi$, and $g(\theta)=\sum_{n\in S}c_n\cos(n\theta)$.
Suppose that
\begin{enumerate}
\item $g\ge0$ on $[-\theta_0,\theta_0]$, with zero set there exactly $\set{\pm\theta_0,\dots,\pm\theta_k}$;
\item some $\lambda_0,\dots,\lambda_k>0$ satisfy $\sum_j\lambda_j\cos(n\theta_j)=0$ for every $n\in S$;
\item the matrix $(\cos(n\theta_j))_{n\in S,\,1\le j\le k}$ has rank $k$, and $(\sin(n\theta_j))_{n\in S,\,0\le j\le k}$ is invertible;
\item some $n_0\in S$ has $\cos(n_0\theta_j)>0$ for $1\le j\le k$.
\end{enumerate}
Then $M(S)=\theta_0/\pi$, and no polynomial with spectrum in $\pm S$ is zero-free on an open interval of length
$\theta_0/\pi$. On the closed interval $I$ of that length centred at $0$, the polynomials with spectrum in $\pm S$ that are
nonnegative on $I$ are the nonnegative multiples of $g(2\pi t)$, and the only probability measure on $I$ annihilating
the spectrum puts mass $\lambda_j/(2\sum_i\lambda_i)$ at each of $\pm\theta_j/(2\pi)$; it has $2k+2$ atoms.
\end{lemma}

\begin{proof}
Write $\theta=2\pi t$. The measure $\mu$ described is a probability measure on $I$; its cosine moments vanish by
(ii) and its sine moments by symmetry. If a polynomial with spectrum in $\pm S$ had no zero on an open interval of
length $\theta_0/\pi$, then after translation and a change of sign it would be strictly positive on the interior of
$I$ and nonnegative at its endpoints, and its integral against $\mu$ would be positive, since $\mu$ charges the
interior points $\pm\theta_j$, $j\ge1$. Larger intervals contain such an interval. Conversely, for $\theta'<\theta_0$ and
small $\delta>0$ the polynomial $g+\delta\cos(n_0\theta)$ is strictly positive on $[-\theta',\theta']$: near the zeros
$\pm\theta_j$, $j\ge1$, by (iv), and elsewhere because $g$ has a positive minimum there. So $M(S)=\theta_0/\pi$.

A polynomial $P$ nonnegative on $I$ integrates to zero against $\mu$, so it vanishes at the $2k+2$ atoms. Its odd
part vanishes at $\theta_0,\dots,\theta_k$, hence is zero by the sine part of (iii). Its even part lies in the left kernel
of $C=(\cos(n\theta_j))_{n\in S,\,0\le j\le k}$; this square matrix is singular by (ii) and has rank at least $k$ by (iii),
so its left kernel is the line through the coefficients of $g$, and the multiple is nonnegative because $g>0$
somewhere on $I$. An annihilating probability measure $\nu$ on $I$ integrates $g$ to zero, so it lives on the atoms of
$\mu$. The differences of its masses at $\pm\theta_j$ satisfy the homogeneous sine system, hence vanish, and the pair
totals lie in the right kernel of $C$, the line through $(\lambda_j)$; normalization gives $\nu=\mu$.
\end{proof}

For $\eps\ge0$ put $h_\eps(s)=\sin(\eps s)/\eps$ if $\eps>0$, $h_0(s)=s$, and $L_\eps=\pi/(2\eps)$, $L_0=+\infty$. Consider
the families
\begin{align*}
 \textup{(T)}\quad &F_\eps=h_\eps+A\cos((1+\eps)s)+B\sin((2+\eps)s), && 0\le\eps\le\tfrac13,\\
 \textup{(H$_\sigma$)}\quad &F_\eps=h_\eps+A\bigl(\cos((1+\eps)s)+\sigma\sin((1+\eps)s)\bigr)+B\cos((2+\eps)s), && 0\le\eps\le\eps_\sigma,\\
 \textup{(Q)}\quad &F_\eps=h_\eps+A\cos((1+\eps)s)+B\sin((2+\eps)s)+C\cos((3+\eps)s), && 0\le\eps\le\tfrac13,
\end{align*}
with $\sigma\in\set{\pm1}$, $\eps_{-1}=2/5$ and $\eps_{+1}=2/7$. In (T) and (H$_\sigma$) the unknowns are $(A,B,u,v,x)$ with
$x<0<u<v<L_\eps$ and the five equations $F_\eps(u)=F_\eps(v)=F_\eps'(u)=F_\eps'(v)=F_\eps(x)=0$; in (Q) they are
$(A,B,C,u,v,w,x)$ with $x<0<u<v<w<L_\eps$ and the seven equations $F_\eps=0$ at $u,v,w,x$ and $F_\eps'=0$ at $u,v,w$. Primes
denote $\dd/\dd s$. The contacts other than $x$ are called inner, and the terms of $F_\eps$ other than $h_\eps$ oscillating.

\begin{lemma}[Interval certificates; \stU]\label{zfg:lem:certs}
For each family and each $\eps$ in its range there is a solution with $-2<x<0$, all contacts below $L_\eps$ (and
$v<4$ in (T), $w<5$ in (Q)) and nonsingular Jacobian, such that:
\begin{enumerate}
\item $F_\eps\ge0$ on $[x,L_\eps]$ with zero set exactly the contacts, $F_\eps'(x)>0$, and $F_\eps''>0$ at the inner contacts;
\item the vectors of values of the oscillating terms at the contacts satisfy a linear relation with positive
coefficients $\lambda_s$, and those at the inner contacts are linearly independent;
\item the square matrix with rows $(\cos(\eps s),\sin((1+\eps)s),\cos((2+\eps)s))$ in (T),
$(\cos(\eps s),\sin((1+\eps)s)-\sigma\cos((1+\eps)s),\sin((2+\eps)s))$ in (H$_\sigma$), and
$(\cos(\eps s),\sin((1+\eps)s),\cos((2+\eps)s),\sin((3+\eps)s))$ in (Q), $s$ running over the contacts, is invertible.
\end{enumerate}
\end{lemma}

\begin{proof}
This is a finite computation in exact rational interval arithmetic. The parameter range is covered by rational
cells $E$ (588 cells for (T), 354 for (H$_{-1}$), 780 for (H$_{+1}$), 1127 for (Q)). For each cell the data give a
rational centre $z_0$, a radius $r$ and a rational preconditioner $R$, which are proposals and not assumptions. With
$g$ the vector of equations and $Z=z_0+[-r,r]^d$, the checker encloses $Rg(E,z_0)$ and $I-RD_zg(E,Z)$ and verifies, for
every row $i$, $q_i=\sum_j\sup\abs{(I-RD_zg)_{ij}}<1$ and $\sup_E\abs{(Rg(\eps,z_0))_i}+rq_i<r$. Then $z\mapsto z-Rg(\eps,z)$ maps $Z$
into its interior and contracts, $R$ and $D_zg$ are invertible, and the fixed point solves $g=0$. On each root box
the checker verifies the ordering, $F'>0$ on a neighbourhood of the enclosure of $x$ and $F''>0$ on neighbourhoods of
the inner contacts, and $F>0$ by bisection in between and up to $\min(c_0,L_\eps)$, where $c_0=7$ in (T) and $9$ in
(H$_\sigma$) and (Q) (using $\pi<22/7$). Beyond $c_0$, concavity of the sine gives $h_\eps(s)\ge2s/\pi\ge2c_0/\pi$, which is
checked to exceed the sum of the absolute values of the oscillating coefficients (weighted by $\sqrt2$ for the first
term of (H$_\sigma$)). The relation (ii) is verified through signs: in (T), $\cos((1+\eps)s)>0$ at $x,u$ and $<0$ at $v$, and
the ratios $\rho=\sin((2+\eps)s)/\cos((1+\eps)s)$ satisfy $\rho(x)<\rho(v)<\rho(u)$; in (H$_\sigma$) the three cyclic $2\times2$
determinants have the same sign; in (Q) the four signed maximal minors $(-1)^j\det V_{\widehat j}$ are negative. The
determinants in (iii) are enclosed away from $0$. All operations round outward on a fixed grid ($10^{-23}$, or
$10^{-24}$ in (Q)), with sine and cosine evaluated by Taylor polynomials at a reduced argument followed by doubling
steps, and with explicit remainder bounds. The data and checkers are
\begin{center}\small
\artifact{kozma_oravecz_three_frequency_2026-11-28/cells.json}, \artifact{kozma_oravecz_three_frequency_2026-11-28/all_odd.js},\\
\artifact{kozma_oravecz_half_integer_2026-11-28/minus.json}, \artifact{kozma_oravecz_half_integer_2026-11-28/plus.json},\\
\artifact{kozma_oravecz_half_integer_2026-11-28/certify_family.js},\\
\artifact{proveit_three_extensions_2026-11-27/checks/spectral/four_frequency_cells.json},\\
\artifact{proveit_three_extensions_2026-11-27/checks/spectral/four_frequency_certificate.py}.
\end{center}
\end{proof}

Within each family the certified solution is unique at each parameter. Its oscillating terms and $h_\eps$ vanish
against the positive weights $\lambda_s$ of (ii), since $F_\eps$ vanishes at the contacts. If two certified solutions had
endpoints $x_1<x_2$, the first $F_\eps$ would be nonnegative at the contacts of the second and positive at $x_2$, so its
sum against the second weights would be positive instead of zero. For equal endpoints the same sum forces the
contact sets to agree, and the independence in (ii) fixes the coefficients. With the nonsingular Jacobian, the
analytic implicit function theorem gives analytic branches on the cells that agree at common endpoints; hence each
family has a real analytic branch on a neighbourhood of its parameter range, in particular at $\eps=0$. We write
$x(\eps)$ for (T), $x_\sigma(\eps)$ for (H$_\sigma$) and $x^{(4)}(\eps)$ for (Q).

\begin{theorem}[Three frequencies, odd-integer step ratio; \stU]\label{zfg:thm:three}
Let $N\ge1$ and let $m\ge3$ be odd. Then
\[
 N\,M(\set{N,(m+1)N,(2m+1)N})=\frac12-\frac{x(1/m)}{\pi m}.
\]
The supremum is not attained, even on an open interval. On a closed interval of the critical length the
nonnegative polynomials with this spectrum form a single ray, and exactly one probability measure annihilates
the spectrum; it has six atoms. As $m\to\infty$ through odd values,
\[
 N\,M(\set{N,(m+1)N,(2m+1)N})=\frac12+\frac{c_1}m+\frac{c_2}{m^2}+O(m^{-3}),\qquad c_1=-\frac{x(0)}\pi,\quad c_2=-\frac{x'(0)}\pi,
\]
with $-1.355137<x(0)<-1.355132$ and $0.7859003465<x'(0)<0.7859187030$, so $c_1\approx0.4313527$ and $c_2\approx-0.2501628$.
\end{theorem}

\begin{proof}
By \cref{zfg:prop:duality} take $N=1$. Put $\eps=1/m$ and $s=m(\pi/2-\theta)$. Since $m$ is odd,
\[
 \cos\theta=\eps h_\eps(s),\quad \cos((m+1)\theta)=(-1)^{(m+1)/2}\cos((1+\eps)s),\quad \cos((2m+1)\theta)=-\sin((2+\eps)s),
\]
so $F_\eps$ of (T) is an even polynomial $g(\theta)$ with spectrum $\set{1,m+1,2m+1}$, and $s\in[x,L_\eps]$ corresponds to
$\theta\in[0,\theta_0]$ with $\theta_0=\pi/2-x/m<\pi$. Under this substitution the rows in (iii) of \cref{zfg:lem:certs} are, up to
signs, the sine modes $\sin\theta,\sin((m+1)\theta),\sin((2m+1)\theta)$. Thus (i)--(iii) of \cref{zfg:lem:certs} give (i)--(iii) of
\cref{zfg:lem:cert} with the nodes $\theta_s=\pi/2-s/m$, and (iv) holds with $n_0=1$ because the inner contacts lie in
$(0,L_\eps)$. \Cref{zfg:lem:cert} gives everything except the expansion, which is the Taylor expansion of the analytic
branch at $0$. The derivative $x'(0)$ is obtained by the linear solve $D_zg\,z'(0)=-\partial_\eps g$ at $\eps=0$, and the
enclosures of $x(0)$ and $x'(0)$ come from a separate interval computation
(\artifact{kozma_oravecz_three_frequency_2026-11-28/limit.js}).
\end{proof}

\begin{theorem}[Three frequencies, half-odd-integer step ratio; \stU]\label{zfg:thm:half}
Let $n\ge1$, let $d\ge5$ be odd, and let $\sigma=-1$ if $d\equiv1\pmod4$ and $\sigma=1$ if $d\equiv3\pmod4$. Then
\[
 M(\set{2n,(d+2)n,(2d+2)n})=\frac1{4n}-\frac{x_\sigma(2/d)}{\pi dn}.
\]
The supremum is not attained, even on an open interval; on a closed interval of the critical length there is one ray
of nonnegative polynomials and one annihilating probability measure, with six atoms. With $N=2n$ and $r=d/2$, for
each residue class of $d$ modulo $4$,
\[
 N\,M=\frac12+\frac{c_{1,\sigma}}r+\frac{c_{2,\sigma}}{r^2}+O(r^{-3}),\qquad c_{1,\sigma}=-\frac{x_\sigma(0)}\pi,\quad c_{2,\sigma}=-\frac{x_\sigma'(0)}\pi,
\]
where $c_{1,-1}\approx0.6064836353$, $c_{2,-1}\approx-0.3704115178$, $c_{1,1}\approx0.2224839370$ and $c_{2,1}\approx-0.1254621878$.
\end{theorem}

\begin{proof}
By \cref{zfg:prop:duality} take $n=1$, so the spectrum is $S=\set{2,d+2,2d+2}$. Put $r=d/2$, $\eps=1/r$, $\theta=2\phi$ and
$s=r(\pi/2-\theta)$, where $\phi=2\pi t$. Then $\cos(2\phi)=\eps h_\eps(s)$; since $\tan((1+r)\pi/2)=\sigma$, $\cos((d+2)\phi)=\cos((1+r)\theta)$ is a
nonzero multiple of $\cos((1+\eps)s)+\sigma\sin((1+\eps)s)$; and since $1+2r$ is even, $\cos((2d+2)\phi)$ is $\pm\cos((2+\eps)s)$. The
range $d\ge5$ gives $\eps\le2/5$ for $\sigma=-1$ and $\eps\le2/7$ for $\sigma=1$. As in the proof of \cref{zfg:thm:three},
\cref{zfg:lem:certs} for (H$_\sigma$) gives the hypotheses of \cref{zfg:lem:cert} at the angles $\phi_s=(\pi/2-s/r)/2$, with
$n_0=2$; the rows in (iii) are, up to nonzero factors, the three sine modes. The critical half-angle is
$\phi_0=(\pi/2-x_\sigma(\eps)/r)/2$, and $M=\phi_0/\pi$. The expansion is that of the analytic branch.
\end{proof}

\begin{theorem}[Four frequencies, odd-integer step ratio; \stU]\label{zfg:thm:four}
Let $N\ge1$ and let $m\ge3$ be odd. Then
\[
 N\,M(\set{N,(m+1)N,(2m+1)N,(3m+1)N})=\frac12-\frac{x^{(4)}(1/m)}{\pi m}.
\]
The supremum is not attained, even on an open interval. On the centred closed interval of the critical length the
nonnegative polynomials with this spectrum form one ray, and exactly one probability measure annihilates the spectrum;
it has eight atoms, and every annihilating probability measure of least support diameter is a translate of it. The
branch $x^{(4)}$ is real analytic on a neighbourhood of $[0,1/3]$, and
\[
 N\,M=\frac12+\frac{c_1^{(4)}}m+\frac{c_2^{(4)}}{m^2}+O(m^{-3}),
\]
where
\[
 c_1^{(4)}=-x^{(4)}(0)/\pi=0.45356051377648\ldots,\qquad c_2^{(4)}=-x^{(4)\prime}(0)/\pi=-0.18321967877208\ldots.
\]
In particular $0.63147839169477181<M(\set{1,4,7,10})<0.63147839169477182$.
\end{theorem}

\begin{proof}
With $N=1$, $\eps=1/m$ and $s=m(\pi/2-\theta)$, the four terms of (Q) are nonzero multiples of
$\cos\theta,\cos((m+1)\theta),\cos((2m+1)\theta),\cos((3m+1)\theta)$:
\begin{gather*}
 h_\eps(s)=m\cos\theta,\qquad \cos((1+\eps)s)=(-1)^{\frac{m+1}2}\cos((m+1)\theta),\\
 \sin((2+\eps)s)=-\cos((2m+1)\theta),\qquad \cos((3+\eps)s)=(-1)^{\frac{3m+1}2}\cos((3m+1)\theta).
\end{gather*}
Apply \cref{zfg:lem:certs} and \cref{zfg:lem:cert} as before, with $n_0=1$; a measure of least support diameter lives on an
interval of the critical length, so it is a translate of the one in \cref{zfg:lem:cert}. The decimals are evaluations
of the analytic branch, enclosed separately; the enclosure at $m=3$ is also checked by an independent certificate for
$T_1+aT_4+bT_7+cT_{10}$ in the coordinate $y=\cos\theta$, with integer Chebyshev arithmetic and no trigonometric evaluation
(\artifact{proveit_three_extensions_2026-11-27/checks/spectral/verify_m3_polynomial.py}).
\end{proof}

When $m$ is even, $b/(2N)=m/2$ is an integer, and \cref{rec:thm:wide} gives the value $1/(2N)$ for both spectra; all
frequencies are then odd, and the lowest cosine is zero-free on an open interval of length $1/(2N)$. Approximate
values of $NM$ for the three-frequency spectra are $0.618557753484$, $0.576776857359$, $0.556692411116$ and
$0.544918838243$ at $m=3,5,7,9$, and of $nM(\set{2n,(d+2)n,(2d+2)n})$ they are $0.346549538462$, $0.276881103951$,
$0.309113910011$ and $0.268194148020$ at $d=5,7,9,11$. The leading coefficients $c_{1,-1}\ne c_{1,1}$ show that no single
expansion in $1/r$ describes both residue classes of $d$.

\begin{example}[The spectrum $\set{1,4,7}$; \stU]\label{zfg:ex:147}
Let $s$ be the root within $10^{-18}$ of $1.12268200790234930920$ of
\begin{align*}
 H(s)={}&165888s^{18}-912384s^{16}+1974528s^{14}-2154240s^{12}+1297800s^{10}\\
 &-467124s^8+105672s^6-15092s^4+1274s^2-49,
\end{align*}
and define
\begin{align*}
 3924844\,p={}&-5579891712s^{16}+30121984512s^{14}-63425084544s^{12}+66376108800s^{10}\\
 &-37602505764s^8+12510048822s^6-2542346352s^4+303526195s^2-18986709,\\
 r={}&3s^2-2p-\tfrac74,\qquad e=\frac{2s^3+2sp-2sr}{8p^2-1},\\
 q(y)={}&y^2-sy+p,\qquad C(y)=y^3+2sy^2+ry+e.
\end{align*}
Then $M(\set{1,4,7})=\arccos(a)/\pi$, where $a\approx-0.36390801664214987744$ is a root of $C$, and the other contacts are the
roots $u\approx0.268082359$ and $v\approx0.854599649$ of $q$; thus
\[
 0.61855775348402062<M(\set{1,4,7})<0.61855775348402063 .
\]
The number $2\cos(\pi M(\set{1,4,7}))=2a$ has degree $54$ over $\Q$, with primitive integral minimal polynomial of leading
coefficient $729$, so $M(\set{1,4,7})$ is transcendental.
\end{example}

\begin{proof}
This is the case $m=3$ of \cref{zfg:thm:three}, with an algebraic certificate. The decisive identity is
\[
 64\,q(y)^2C(y)=A\,T_1(y)+B\,T_4(y)+T_7(y),\qquad A\approx1.6408563,\quad B\approx1.0859836,
\]
verified by expansion modulo $H(s)$, with $B=64p^2e$ from the constant coefficient and $A=7+[y]\,(64q^2C)$ from the linear
one; $8p^2-1$ is invertible modulo $H$. The cubic $C$ increases on $[a,1]$, since $C'(a)>0$ and $C''(y)=6y+4s>0$ there, so
the left side is nonnegative on $[a,1]$ with contacts exactly $a,u,v$. In the projected coordinates $(y,T_4(y))$ the
determinants of consecutive columns at $(u,v)$, $(v,a)$, $(a,u)$ lie in $(-0.554,-0.552)$, $(-0.141,-0.140)$ and
$(-0.192,-0.191)$, which gives positive weights annihilating $T_1$ and $T_4$, and the identity then annihilates $T_7$. The
sine matrix $(U_{n-1}(z))_{n\in\set{1,4,7},\,z\in\set{a,u,v}}$, $U_{n-1}(\cos t)=\sin(nt)/\sin t$, has positive determinant by rational
bounds. So \cref{zfg:lem:cert} applies with $\theta_0=\arccos a$, $\theta_1=\arccos u$, $\theta_2=\arccos v$.

For the degree, work in $\Q[s]/(H)$ and take the primitive integral part with positive leading coefficient of
$\Res_s(H(s),C(y/2))$. It is a polynomial $Q$ of degree $54$ with leading coefficient $729$ and a root at $y=2a$. Its monic
reduction $Q_0$ modulo $211$ satisfies $y^{211^{54}}\equiv y\pmod{Q_0}$ and $\gcd(Q_0,y^{211^i}-y)=1$ for $i=18,27$; since the
prime divisors of $54$ are $2$ and $3$, $Q_0$ is irreducible over $\mathbb F_{211}$ and $Q$ over $\Q$. As $Q$ is primitive and
not monic, $2a$ is not an algebraic integer, whereas a rational $M$ would make $2a=\zeta+\zeta^{-1}$ for a root of unity
$\zeta$. \Cref{zfg:cor:dichotomy} gives transcendence. The checks are in
\artifact{positive-fourier-zeros/kozma/check_gaps.py} and \artifact{positive-fourier-zeros/kozma/verification.txt}.
\end{proof}

The polynomial was found from the normal $D$ of \cref{zfg:thm:threerule}: equating the ratios of its coordinates at $u,v$
and dividing the antisymmetric minors by $u-v$ gives symmetric equations in $s=u+v$, $p=uv$; elimination yields $H$
among the possible factors, the linear subresultant gives $p$, and the double contacts force the factor $q^2$. The sign
checks above identify this branch as the global optimum without locating the real zeros of the other factors.

\begin{remark}[Independent certificates]\label{zfg:rem:147}
In the normalization $P(y)=T_4(y)+Ay+B_*T_7(y)$ the same optimum has $B_*\in[0.9208242456520,0.9208242456522]$ and double
contacts $u\in[0.2680823592348,0.2680823592351]$, $v\in[0.8545996486673,0.8545996486676]$, isolated by monotone root
branches of $\partial_y(P/y)$; $P=(y-u)^2(y-v)^2G$ with a cubic $G$ increasing on $[a,1]$, and
\[
 -0.36390801664214987744459<a<-0.36390801664214987744456,
\]
checked with outward-rounded integer intervals and exact rational endpoints by
\begin{center}\artifact{kozma_oravecz_147_2026-11-27/certify.js} and \artifact{kozma_oravecz_147_2026-11-27/verify.py}.\end{center}
For $d=5$ in
\cref{zfg:thm:half}, the polynomial $P=T_2+A_0T_7+B_0T_{12}$ with
\[
 B_0\in[0.46688711936350,\,0.46688711936362]
\]
has double contacts $u_0\approx0.9422828922$ and $v_0\approx0.7295482273$ and a simple contact $\xi\in(0.46362209927,0.46362209930)$; after dividing out
$(y-\xi)(y-u_0)^2(y-v_0)^2$, the quotient of degree seven has $256$ positive Bernstein bounds on $32$ subintervals of
$[0.46362209927,1]$, and $M(\set{2,7,12})=\arccos(\xi)/\pi\approx0.346549538462$, checked by
\artifact{kozma_oravecz_half_integer_2026-11-28/certify_2712.js}.
\end{remark}

\subsection{Square frequencies}\label{zfg:sec:squares}
\sources{ProveIt_New_Results_2026-11-27; positive-fourier-zeros}

For integers $N\ge1$ and $K\ge0$ put $H=K+1$, $M_{N,K}=M(\set{N^2,(N+1)^2,\dots,(N+K)^2})$ and
$M_{N,\infty}=\sup_{K\ge0}M_{N,K}$; $M_{N,K}$ is nondecreasing in $K$. Kozma and Oravecz showed $M_{N,K}=O(1/N)$ uniformly in $K$
and asked whether this order is correct (Question~2 of \cite{KozmaOravecz2003Gaps}). It is, exactly when $K$ is at least a
fixed multiple of $N$.

\begin{theorem}[Uniform order for squares; \stU]\label{zfg:thm:squares}
For all $N\ge1$ and $K\ge0$,
\[
 \max_{1\le m\le H}\frac{m^2}{2\sum_{j=0}^{m-1}(N+j)^2}\ \le\ M_{N,K}\ \le\ \frac12\sum_{j=0}^{H-1}\frac1{(N+j)^2},
\]
and the left side is attained at the unique prefix length
\[
 m_*=\min(H,m_N),\qquad m_N=\floor*{\frac{1+\sqrt{12N^2-12N+1}}2}.
\]
Moreover
\[
 \frac{3\min(H,N)}{14N^2}\ \le\ M_{N,K}\ \le\ \frac H{N(N+H)},
\]
so $M_{N,K}\asymp\min(K+1,N)/N^2$ with absolute constants.
\end{theorem}

\begin{proof}
The first display is \cref{zfg:prop:duality}. Write the prefix bound as $L_N(m)=m/(A'+B'm+C'm^2)$ with
$A'=2N^2-2N+\tfrac13$, $B'=2N-1$, $C'=\tfrac23$. Cross multiplication gives $L_N(m+1)>L_N(m)$ if and only if
$A'>C'm(m+1)$, that is $m(m+1)\le3N(N-1)$; there are no ties, since the unreduced threshold is
$m(m+1)<3N(N-1)+\tfrac12$. The sign changes once, which gives $m_*$. For the lower estimate take $m=\min(H,N)$: the
average of the first $N$ squares of the band is $(14N^2-9N+1)/6\le\tfrac73N^2$. For the upper estimate use
$n^{-2}\le2/(n(n+1))$ and telescope.
\end{proof}

Since $\min(H,N)/N^2\ge H/(N(N+H))\ge\min(H,N)/(2N^2)$, the two sides differ by at most the factor $14/3$. Short bands are
determined to fourth order.

\begin{theorem}[Short bands; \stU]\label{zfg:thm:shortsq}
Suppose $(N+K)^2\le3N^2$, and put $C=N+K/2$ and $V=K(K+2)/12$. Then
\[
 1\ \le\ \frac{M_{N,K}}{H/(2(C^2+V))}\ \le\ 1+\frac{K(K-1)(K+2)(K+3)}{180\,C^2(C^2+V)}.
\]
In particular $M_{N,K}=\frac{H}{2(C^2+V)}\bigl(1+O((K/N)^4)\bigr)$ for $K=o(N)$, and the cases $K=0,1$ are exact.
\end{theorem}

\begin{proof}
For $a<b\le3a$ the three-point rule with masses $1/(2(1-c))$ at $\pm1/(a+b)$ and $-c/(1-c)$ at $0$, where
$c=\cos(2\pi a/(a+b))\le0$, cancels both frequencies on a span $2/(a+b)$ in circle units (this is \cref{fil:cor:pair} with
$m=1$). Pair the indices $j$ and $K-j$ of the band, using the two-point rule at a middle index. The pair means are
$C^2+(j-K/2)^2$. With $Y=(j-K/2)^2$ for $j$ uniform on $\set{0,\dots,K}$, convolution and the prefix bound with $m=H$ give
\[
 \frac H{2(C^2+\mathbb EY)}\le M_{N,K}\le\frac H2\,\mathbb E\frac1{C^2+Y}.
\]
Here $\mathbb EY=V$ and $\operatorname{Var}Y=K(K-1)(K+2)(K+3)/180$, and with $A=C^2+V$ the identity
$A\,\mathbb E[(C^2+Y)^{-1}]-1=\mathbb E[(Y-V)^2/(A(C^2+Y))]$ gives the bound. The fourth-order gain comes from the cancellation of
the terms linear in $j-K/2$ in the pair means.
\end{proof}

For example $0.0004984141368<M_{100,10}<0.0004984144560$, outward roundings of the two sums in the proof.

\begin{theorem}[Long bands; \stU]\label{zfg:thm:longsq}
As $N\to\infty$:
\begin{enumerate}
\item if $H/N\to\tau\in(0,\infty)$ and $u=\min(\tau,\sqrt3)$, then
\[
 \frac{3u}{2(u^2+3u+3)}\le\liminf NM_{N,K}\le\limsup NM_{N,K}\le\frac\tau{2(1+\tau)};
\]
\item $NM_{N,\infty}\le\tfrac{21}{62}+o(1)$;
\item if $K/N\to\infty$, then $\sqrt3-\tfrac32\le\liminf NM_{N,K}\le\limsup NM_{N,K}\le\tfrac{21}{62}$, and the same holds for $NM_{N,\infty}$.
\end{enumerate}
\end{theorem}

The constants in (iii) are about $0.23205$ and $0.33871$; whether $NM_{N,\infty}$ converges is open. The proof of (ii) uses the
fact that the equal atoms at $\pm1/4$ cancel the frequencies $1,3,7$ together.

\begin{lemma}[Stable cancellation near $1:3:7$; \stU]\label{zfg:lem:stable137}
For every $\eta>0$ all triples $(1,a,b)$ sufficiently close to $(1,3,7)$ admit a probability measure on an interval of
length $\tfrac12+\eta$ whose Fourier transform, in circle units, vanishes at $1,a,b$.
\end{lemma}

\begin{proof}
Convolve $\tfrac12(\delta_{-1/4}+\delta_{1/4})$ with the uniform law on $[-\eta/2,\eta/2]$; the transform keeps the three zeros.
The origin is interior to the convex hull of the six cosine and sine coordinates on the enlarged interval:
otherwise a nonzero combination would be nonnegative there with integral zero against a measure with positive
density on an open interval, contradicting analyticity and the linear independence of the six functions. The
minimum over unit vectors of the maximal scalar product with the curve is then positive, and it moves by at most
the size of a uniform perturbation of the curve, so the origin stays in the hull for nearby frequencies.
\end{proof}

\begin{proof}[Proof of \cref{zfg:thm:longsq}]
For a prefix with $m/N\to t>0$, Riemann sums give $NL_N(m)\to t^2/(2\int_1^{1+t}x^2\dd x)=3t/(2(t^2+3t+3))$, which increases
up to $t=\sqrt3$ and then decreases, with value $\sqrt3-\tfrac32$ at $t=\sqrt3$; the upper sum of \cref{zfg:thm:squares} tends
to $\tfrac12\int_1^{1+\tau}x^{-2}\dd x$. This gives (i) and the lower bound in (iii), by monotonicity in $K$.

For (ii) partition the indices $n\ge N$ into triples with size ratios tending to $1:\sqrt3:\sqrt7$. Choose disjoint
periodic sets $U,V$ of positive integers formed by $58$ and $38$ residue classes modulo $100$, with increasing
enumerations $u_k=k/0.58+O(1)$ and $v_k=k/0.38+O(1)$, and put $P(n)=u_{\floor{0.58\sqrt3\,n}}$ and $Q(n)=v_{\floor{0.38\sqrt7\,n}}$.
Both indexing slopes exceed one, so these maps are injective with disjoint ranges, $P(n)=\sqrt3\,n+O(1)$ and
$Q(n)=\sqrt7\,n+O(1)$, and both exceed $n$ for large $N$. Scanning upwards from $N$, take the triple $(n,P(n),Q(n))$ whenever
$n$ is still unused; its larger members can be neither later starting points nor targets of another start, so this
partitions the tail. By \cref{zfg:lem:stable137}, scaled, the three squares of a triple have a cancelling measure of
diameter at most $(\tfrac12+o(1))/n^2$, uniformly in $n\ge N$, against the sum of the separate two-point widths
$\tfrac1{2n^2}(1+\tfrac13+\tfrac17+O(1/N))$; the ratio is $\tfrac{21}{31}+o(1)$. For a finite band, convolve the laws of the
finitely many triples meeting it. Hence $M_{N,\infty}\le(\tfrac{21}{31}+o(1))\tfrac12\sum_{n\ge N}n^{-2}=(\tfrac{21}{62}+o(1))/N$, which gives
(ii) and, with $M_{N,K}\le M_{N,\infty}$, the upper bound in (iii).
\end{proof}

The partition argument works with any fixed list of distinct odd integers $h_1,\dots,h_l>1$ with $\sum_ih_i^{-1/2}<1$,
using periodic sets of rational densities larger than $h_i^{-1/2}$; it replaces $\tfrac{21}{62}$ by
$[2(1+\sum_ih_i^{-1})]^{-1}$, and $3,7,625$ already give a slightly smaller constant.

\subsection{Higher-power frequencies}\label{zfg:sec:powers}
\sources{ProveIt_New_Results_2026-11-27}

The prefix and convolution bounds of \cref{zfg:prop:duality} apply equally to higher powers.

\begin{corollary}[Power-frequency bands; \stU]\label{zfg:cor:powers}
Fix an integer $d\ge2$, and put $M^{(d)}_{N,H}=M(\set{N^d,(N+1)^d,\dots,(N+H-1)^d})$ and $M_N^{(d),\infty}=\sup_{H\ge1}M^{(d)}_{N,H}$. Then
$M^{(d)}_{N,H}\asymp_d\min(H,N)/N^d$ uniformly in $N,H\ge1$, and $M^{(d)}_{N,H}=H(1+O_d(H/N))/(2N^d)$ if $H=o(N)$. Moreover
\[
 \sup_{t>0}\frac{(d+1)t^2}{2((1+t)^{d+1}-1)}\ \le\ \liminf N^{d-1}M_N^{(d),\infty}\ \le\ \limsup N^{d-1}M_N^{(d),\infty}\ \le\ \frac1{2(d-1)},
\]
and the supremum on the left is attained at the unique positive solution $t_d$ of $(1+t)^d(2-(d-1)t)=2$.
\end{corollary}

\begin{proof}
Apply \cref{zfg:prop:duality} with $a_j=(N+j-1)^d$. The prefix $m=\min(H,N)$ gives the lower bound $m/(2(2N)^d)$; the upper
bound is at most $H/(2N^d)$ and at most $\tfrac12\sum_{n\ge N}n^{-d}=O_d(N^{1-d})$. In a short band $(N+j)^d=N^d(1+O_d(H/N))$
uniformly. For $m/N\to t$ the scaled prefix bound tends to $t^2/(2\int_1^{1+t}x^d\dd x)$, and integral comparison gives the
upper constant. The derivative of $(1+t)^d(2-(d-1)t)$ is $(d+1)(1+t)^{d-1}(1-(d-1)t)$, so this function starts at $2$, increases
until $t=1/(d-1)$ and decreases to $0$ at $t=2/(d-1)$; it crosses $2$ exactly once, at the maximizer.
\end{proof}

\subsection{Reciprocal moments of zero heights and support uncertainty}\label{zfg:sec:recip}
\sources{ProveIt_New_Results_2026-11-27; positive-fourier-zeros}

Let $\mu$ be a compactly supported probability measure on $\R$ with $L=\zfgdiam\supp\mu>0$, and $G(z)=\int e^{zt}\dd\mu(t)$. A
selected multiset of zeros of $G$ uses each zero at most as often as its multiplicity; distinct zeros may have the
same imaginary part. \Cref{fil:thm:meanbound} bounds the mean height of zeros on one vertical line; the same bound holds
for the prefixes of the heights of arbitrary complex zeros.

\begin{proposition}[Prefix bound; \stU]\label{zfg:prop:prefix}
Let $\rho_1,\dots,\rho_n$ be selected zeros of $G$ with heights $0<\lambda_1\le\dots\le\lambda_n$, $\lambda_j=\Im\rho_j$. Then
\[
 S_m=\sum_{j=1}^m\lambda_j\ \ge\ \frac{\pi m^2}L\qquad(1\le m\le n).
\]
In particular at most $LT/\pi$ selected zeros have height at most $T$, so the heights of the zeros of $G$ in the upper
half-plane do not accumulate at any finite height.
\end{proposition}

\begin{proof}
Fix $m$. For each distinct selected zero $\rho$ of selected order $k$ take the real and imaginary parts of $t^\ell e^{\rho t}$,
$0\le\ell<k$. Their integral against $\mu$ vanishes and lies in the convex hull of their values on $\supp\mu$, so by
Carath\'eodory's theorem an atomic probability measure on the same interval keeps every selected zero to its order.
Translate its leftmost atom to $0$ and let $D\le L$ be its diameter; $D>0$, since a single atom has no zeros. Put
$\kappa=S_m/m$ and apply \cref{fil:lem:strip} with $T=2\kappa$: the selected zeros contribute at least
$\sum_{j\le m}(T-\lambda_j)_+\ge mT-S_m=m\kappa$, other zeros only increase the sum, so $0\le J(T)\le2D\kappa^2-2\pi m\kappa$, that is
$S_m\ge\pi m^2/D\ge\pi m^2/L$. If $m$ selected heights are at most $T$, then $\pi m^2/L\le mT$.
\end{proof}

\begin{theorem}[Sharp reciprocal moments of zero heights; \stU]\label{zfg:thm:recip}
Let $\rho_1,\dots,\rho_n$ be selected zeros of $G$ with heights $0<\lambda_1\le\dots\le\lambda_n$. For every real $p>0$,
\[
 \sum_{j=1}^n\lambda_j^{-p}\ \le\ \Bigl(\frac L\pi\Bigr)^p\sum_{j=1}^n\frac1{(2j-1)^p}.
\]
Equality holds if and only if, for some $u\in\R$ and $0<b<1$,
\[
 \mu=b\delta_u+(1-b)\delta_{u+L},\qquad \rho_j=\frac1L\log\frac b{1-b}+\frac{(2j-1)\pi i}L\quad(1\le j\le n);
\]
if all $\rho_j$ lie on the imaginary axis, equality requires $b=\tfrac12$. For an arbitrary finite or infinite selected
multiset and $p>1$,
\[
 \sum_j\lambda_j^{-p}\ \le\ \Bigl(\frac L\pi\Bigr)^p(1-2^{-p})\zeta(p),
\]
with equality exactly for the measures above with all $j\ge1$. In particular $\sum_j(\Im\rho_j)^{-2}\le L^2/8$.
\end{theorem}

\begin{proof}
Let $c=\pi/L$ and $v_j=c(2j-1)$, and put $H_\lambda(t)=\sum_{j\le n}(t-\lambda_j)_+$ and $H_v(t)=\sum_{j\le n}(t-v_j)_+$. For a sorted list,
$H_\lambda(t)=\max_{0\le m\le n}(mt-S_m)$, the maximizing prefix consisting of the entries below $t$. Since $\sum_{j\le m}v_j=cm^2$,
\cref{zfg:prop:prefix} gives $H_\lambda\le H_v$ everywhere. For $x,p>0$,
\[
 x^{-p}=p(p+1)\int_0^\infty(t-x)_+t^{-p-2}\dd t,
\]
and integrating $H_\lambda\le H_v$ against this positive weight gives the finite inequality. At equality the continuous
function $H_v-H_\lambda\ge0$ has integral zero against a positive weight, so it vanishes; its derivative jumps recover the
multisets, so $\lambda_j=(2j-1)\pi/L$. Then there is a zero $\rho=\beta+i\pi/L$. With the support translated to $[0,L]$, the
imaginary part of $G(\rho)=0$ is $\int e^{\beta t}\sin(\pi t/L)\dd\mu(t)=0$, whose integrand is positive on $(0,L)$; so $\mu$ lives on
the two endpoints, both charged since the diameter is $L$. With masses $b,1-b$ the real part gives $e^{\beta L}=b/(1-b)$, and the
zeros of $b+(1-b)e^{Lz}$ are those displayed. Conversely these measures attain the bound, and $\beta=0$ forces $b=\tfrac12$.
For an infinite multiset, \cref{zfg:prop:prefix} permits increasing enumeration; apply the finite inequality to every
prefix and let $n\to\infty$. For equality, the hinge comparison holds with a supremum over all $m\ge0$, and Tonelli's
theorem forces equality of the infinite hinge functions, so all odd heights occur and the endpoint argument applies.
Finally $(1-2^{-2})\zeta(2)=\pi^2/8$.
\end{proof}

The heights are bounded without any restriction on the real parts of the zeros; a positive measure need not have
any zero on the imaginary axis. The proof extends to every convex decreasing statistic.

\begin{corollary}[Decreasing convex statistics; \stU]\label{zfg:cor:convex}
Let $f:(0,\infty)\to[0,\infty)$ be twice continuously differentiable, decreasing and convex, with $f(t)\to0$ as $t\to\infty$.
Then $\sum_{j\le n}f(\lambda_j)\le\sum_{j\le n}f((2j-1)\pi/L)$, and the infinite analogue holds whenever its right side is finite.
In particular, for every selected finite or infinite multiset,
\[
 \sum_je^{-s\lambda_j}\le\frac1{2\sinh(\pi s/L)}\quad(s>0),\qquad \prod_j\Bigl(1+\frac{x^2}{\lambda_j^2}\Bigr)\le\cosh\frac{Lx}2\quad(x\ge0),
\]
with equality for some $s>0$ or $x>0$ exactly in the equality case of \cref{zfg:thm:recip}.
\end{corollary}

\begin{proof}
Convexity and decay give $0\le-tf'(t)\le2f(t/2)\to0$, so integration by parts gives $f(x)=\int_0^\infty(t-x)_+f''(t)\dd t$, and the
hinge comparison applies. The first bound sums a geometric series. For the second take $f(u)=\log(1+x^2/u^2)$, with
$f''(u)=2x^2(3u^2+x^2)/(u^2(u^2+x^2)^2)>0$, and use $\cosh y=\prod_{j\ge1}(1+4y^2/((2j-1)^2\pi^2))$; the product converges by
\cref{zfg:thm:recip}. Strict positivity of $f''$ gives the equality statement.
\end{proof}

\paragraph{Steinerberger's support problem.}
On $\R/\Z$ write $\widehat f(a)=\int f(t)e^{-2\pi iat}\dd t$. Steinerberger \cite[Problem~41, \emph{A kind of uncertainty
principle?}]{Steinerberger2026Open} asks whether a nonnegative $f$ with positive integral, support in $(-r,r)$ and
$\widehat f(a)=0$ for all $a$ in a finite set $A\subset\N$ must have $r$ bounded below by a constant multiple of
$(\sum_{a\in A}a^{-2})^{1/2}$.

\begin{theorem}[Sharp support uncertainty; \stU]\label{zfg:thm:support}
Let $A\subset\N$ be finite and nonempty, $n=\abs A$, and let $f\in L^1(\R/\Z)$ be nonnegative with positive integral,
$\supp f\subset(-r,r)$ and $\widehat f(a)=0$ for $a\in A$. Then
\[
 r\ \ge\ C_n\Bigl(\sum_{a\in A}\frac1{a^2}\Bigr)^{1/2},\qquad C_n=\frac1{4\bigl(\sum_{j=1}^n(2j-1)^{-2}\bigr)^{1/2}},
\]
and in particular $r\ge(\pi\sqrt2)^{-1}(\sum_{a\in A}a^{-2})^{1/2}$. Each $C_n$ and the constant $1/(\pi\sqrt2)$ are optimal, also for smooth $f$,
where sharpness is an infimum and is not attained.
\end{theorem}

\begin{proof}
If $r<\tfrac12$, lift $f$ to $(-r,r)\subset\R$ and normalize it to a probability measure of diameter at most $2r$ whose transform
vanishes at the angular frequencies $2\pi a$. The case $p=2$ of \cref{zfg:thm:recip} gives
$\tfrac1{4\pi^2}\sum_{a\in A}a^{-2}\le\tfrac{4r^2}{\pi^2}\sum_{j\le n}(2j-1)^{-2}$. If $r\ge\tfrac12$ the bound is automatic, since its right
side is at most $\pi/(4\sqrt6)<\tfrac12$. Letting the odd sum increase to $\pi^2/8$ gives the uniform constant. For sharpness take
$A=\set{1,3,\dots,2n-1}$ and $\tfrac12(\delta_{-1/4}+\delta_{1/4})$, whose odd Fourier coefficients vanish, and convolve with a smooth
probability density supported in $(-\eta,\eta)$; let $\eta\downarrow0$ and then $n\to\infty$. Equality for absolutely continuous
measures is excluded by the equality case of \cref{zfg:thm:recip}.
\end{proof}

On the circle $\R/(2\pi\Z)$ with characters $e^{-iat}$ the radius scales by $2\pi$, and the optimal uniform constant is $\sqrt2$.
The theorem proves the inequality asked for, with its best constant; it does not determine the least radius for a
given set $A$. Near the extremal configuration the measure is close to the two endpoint atoms, linearly in the
deficit.

\begin{theorem}[Stability of the extremizer; \stU]\label{zfg:thm:stab}
Let $\mu$ be a probability measure on $[-r,r]$ whose transform $\widehat\mu(a)=\int e^{2\pi iat}\dd\mu(t)$ has the selected positive
real zeros $a_1\le\dots\le a_m$, and put
\[
 \mathcal D=\sum_{j=1}^m\Bigl(\frac1{(2j-1)^2}-\frac1{(4ra_j)^2}\Bigr),\qquad \nu_r=\tfrac12(\delta_{-r}+\delta_r).
\]
If $\mathcal D<\tfrac12$, then $0\le4ra_1-1\le2\mathcal D$ and $W_1(\mu,\nu_r)\le(2\pi+\pi^2/2)\,r\,\mathcal D$, uniformly in $m$.
\end{theorem}

\begin{proof}
Put $x_j=4ra_j$, $\delta=x_1-1\ge0$ (by \cref{zfg:prop:prefix} with $L\le2r$), and let $H_x,H_b$ be the hinge functions of
$(x_j)$ and of $b_j=2j-1$. For $1\le t\le2$ and $k\ge2$, $kt-k^2\le0$; hence $H_b(t)=t-1$ and, by the prefix bound,
$H_x(t)=(t-1-\delta)_+$ on $[1,2]$. With $d=\min(\delta,1)$ and the integral representation of $x^{-2}$,
\[
 \mathcal D\ \ge\ 6\int_1^2\min(d,t-1)\,t^{-4}\dd t=1-(1+d)^{-2}-\frac d4\ \ge\ \frac d2,
\]
the last step by concavity on $[0,1]$ and the values at the endpoints. So $\delta\le2\mathcal D$ when $\mathcal D<\tfrac12$.

Let $U=t/r$ under $\mu$. Since $\mathbb Ee^{i\pi x_1U/2}=0$, the Lipschitz bound for the cosine gives $\mathbb E\cos(\pi U/2)\le\pi\delta/2$, and
$\cos(\pi U/2)\ge1-\abs U$ on $[-1,1]$, so $d_0=\mathbb E(1-\abs U)\le\pi\delta/2$. Sending $U$ to its sign $S$ (either endpoint at $0$) costs
$d_0$ in normalized units, and the sine moment gives
$\abs{\mathbb ES}\le\mathbb E\abs{S-\sin(\pi x_1U/2)}\le\tfrac\pi2(d_0+\delta)$. Moving the excess endpoint mass $\abs{\mathbb ES}/2$ across costs
$\abs{\mathbb ES}$. Undoing the scale and using $\delta\le2\mathcal D$ gives the bound.
\end{proof}

With $L=2r$ and the full deficit $\Delta=\pi^2/8-(\pi/L)^2\sum_j(2\pi a_j)^{-2}$, one has $\Delta\ge\mathcal D$ for a finite selected
multiset and $\Delta=\lim_m\mathcal D$ along the prefixes of an infinite one; hence $W_1(\mu,\nu_r)\le(\pi+\pi^2/4)L\Delta$ whenever
$\Delta<\tfrac12$.

\begin{remark}[Higher dimensions]\label{zfg:rem:dim}
Steinerberger also suggests higher dimensions. The direct analogue with $\abs a^{-2}$ fails on $(\R/\Z)^d$ for every $d\ge2$:
the atoms at $\pm e_1/4$ have transform $\cos(\pi a_1/2)$, which vanishes at every lattice vector with odd first
coordinate, and the vectors with $2^k\le a_1<2^{k+1}$ odd and $2^k\le a_2<2^{k+1}$ contribute an amount bounded below at each
scale $k$, so $\sum\abs a^{-2}$ diverges, while a smooth convolution keeps the support in a ball of radius less than $\tfrac12$.
\end{remark}

\subsection{The Lorentz degree of polynomials with zeros in a thin ellipse}\label{zfg:sec:lorentz}
\sources{erdelyi_ellipse_and_proveit_extensions}

A Lorentz representation of a real polynomial $P$ on $[-1,1]$ is an identity
\[
 P(x)=\sum_{j=0}^Da_j(1-x)^j(1+x)^{D-j}
\]
with all $a_j$ of the same weak sign; the least such $D$ is the Lorentz degree $\zfgdL(P)$, infinite if there is none. If $P>0$
on $[-1,1]$, every such representation has nonnegative coefficients. In the variable $z=(1-x)/(1+x)$ a representation of
degree $D$ is the statement that $(1+z)^DP((1-z)/(1+z))$ has nonnegative coefficients, so the Lorentz degree is the least
number of factors $1+z$ that makes a polynomial with prescribed roots coefficientwise nonnegative. Positivity of the same kind decides sharpness in
\cref{fil:thm:reflect}, where a central root of the filter polynomial supplies the same factor. For $0<\eps\le1$ let
\[
 \zfgE_\eps=\set{x+iy:x^2+y^2/\eps^2<1},
\]
and let $L_n(\eps)$ be the least finite Lorentz degree of a real polynomial of degree $n$ with all zeros in $\zfgE_\eps$.
A real polynomial of odd degree with all zeros in $\zfgE_\eps$ has a zero in $(-1,1)$, where a nonzero polynomial with a
same-sign representation cannot vanish; so $n$ is even. Erd\'elyi \cite[\S8, p.~197]{Erdelyi1991Estimates} asked whether
$\zfgdL(p)\ge cn/\eps^2$, with an absolute $c>0$, for every polynomial $p$ of degree $n$ with all zeros in $\zfgE_\eps$.

\begin{theorem}[The order of the least Lorentz degree; \stU]\label{zfg:thm:lorentz}
For every even $n\ge2$ and $0<\eps\le1$, with $M_n(\eps)=\max\set{n/\eps,\eps^{-2}}$,
\[
 \frac3{100}\,M_n(\eps)\ \le\ L_n(\eps)\ \le\ 513\,M_n(\eps).
\]
The upper bound is attained by polynomials with rational coefficients and simple zeros, and for $n\eps>4$ its constant
may be taken to be $64$.
\end{theorem}

The lower bound combines \cref{zfg:thm:ellipselow,zfg:thm:ellipsewidth}, and the upper bound is \cref{zfg:thm:ellipseup}.
The following family already answers Erd\'elyi's question negatively, and defeats every bound $c_\alpha n/\eps^\alpha$ with fixed
$\alpha>1$.

\begin{theorem}[Scaled Chebyshev polynomials; \stU]\label{zfg:thm:erdelyi}
For $n\equiv0\pmod4$ put $s_n=1+n^{-2}$, $P_n(x)=2+T_n(s_nx)$ and $\eps_n=2/n$. Then $P_n\in\Q[x]$ has degree $n$, is strictly positive
on $[-1,1]$, has $n$ distinct zeros in $\zfgE_{\eps_n}$, and
\[
 \frac{s_n^2n^2}3\ \le\ \zfgdL(P_n)\ \le\ 8n^2 .
\]
In particular $\eps_n^2\zfgdL(P_n)/n\le32/n\to0$.
\end{theorem}

The upper bounds rest on the Bernstein operator $B_Df(t)=\sum_{j=0}^Df(j/D)\binom Djt^j(1-t)^{D-j}=\mathbb E[f(K/D)]$, $K\sim\mathrm{Binomial}(D,t)$,
on $[0,1]$. It maps polynomials of degree at most $n\le D$ to themselves, since the falling factorials of $K$ have
expectations $(D)_jt^j$.

\begin{lemma}[Inverting the Bernstein operator; \stU]\label{zfg:lem:bernstein}
On polynomials of degree at most $n\le D$ with the supremum norm on $[0,1]$, $\norm{I-B_D}\le2n^2/D$. If $\eta=2n^2/D<1$, then $B_D$
is invertible there and $\norm{B_D^{-1}-I}\le\eta/(1-\eta)$.
\end{lemma}

\begin{proof}
Let $g(\theta)=f((1+\cos\theta)/2)$ and $t=(1+\cos\theta)/2$; then $g''(\theta)=t(1-t)f''(t)-\tfrac12(2t-1)f'(t)$. Bernstein's inequality for
trigonometric polynomials, applied twice, gives $\norm{g''}\le n^2\norm f$ \cite{QueffelecZarouf2019Bernstein}, and Markov's
inequality on $[0,1]$ gives $\norm{f'}\le2n^2\norm f$, so $H=\sup_{0<t<1}t(1-t)\abs{f''(t)}\le2n^2\norm f$. With
$h(t)=t\log t+(1-t)\log(1-t)$, $h''=1/(t(1-t))$, the functions $Hh\pm f$ are convex, and their tangent inequalities at $t$ give
\[
 \abs{f(u)-f(t)-(u-t)f'(t)}\le H\Bigl[u\log\frac ut+(1-u)\log\frac{1-u}{1-t}\Bigr]\le H\frac{(u-t)^2}{t(1-t)}\qquad(0\le u\le1).
\]
Take $u=K/D$ and expectations, using $\mathbb EK/D=t$ and $\operatorname{Var}(K/D)=t(1-t)/D$; at $t=0,1$ the error vanishes. The
Neumann series gives the inverse.
\end{proof}

\begin{proof}[Proof of \cref{zfg:thm:erdelyi}]
Let $a=\operatorname{arcosh}2$. The zeros are $z_j=s_n^{-1}\cos((2j+1)\pi/n+ia/n)$, $0\le j<n$, since $T_n$ multiplies the argument by $n$
and gives $-\cosh a=-2$; they are distinct points of an ellipse with semiaxes $A_n=\cosh(a/n)/s_n$ and $B_n=\sinh(a/n)/s_n$.
Termwise comparison of the series gives $\cosh(a/n)<1+(\cosh a-1)/n^2=s_n$ and $\sinh(a/n)\le\sinh(a)/n=\sqrt3/n$, so $A_n<1$,
$B_n<2/n$, and all zeros lie in $\zfgE_{2/n}$. On $[-1,1]$, $P_n\ge1$ (for $\abs{s_nx}\ge1$, $T_n(s_nx)\ge1$ as $n$ is even), and
$\norm{T_n(s_n\,\cdot)}\le T_n(s_n)=\cosh(n\operatorname{arcosh}(1+n^{-2}))\le\cosh\sqrt2<\tfrac94$.

Upper bound: with $D=8n^2$, $\eta=\tfrac14$, $F_n(t)=P_n(2t-1)$ and $B_D1=1$, \cref{zfg:lem:bernstein} gives
$\norm{B_D^{-1}F_n-F_n}\le\tfrac13\norm{T_n(s_n\,\cdot)}<\tfrac34$. Since $F_n\ge1$, all values $(B_D^{-1}F_n)(j/D)$ are positive, and
$B_D(B_D^{-1}F_n)=F_n$ is a positive Lorentz representation of degree $D$ in $x=2t-1$.

Lower bound: $h_j(x)=(1-x)^j(1+x)^{D-j}$ has $h_j(0)=1$ and $h_j''(0)=(D-2j)^2-D\ge-D$, so every $P>0$ with a representation of
degree $D$ has $P''(0)\ge-DP(0)$. For $4\mid n$, $T_n(0)=1$ and the Chebyshev equation gives $T_n''(0)=-n^2$, so $P_n(0)=3$ and
$P_n''(0)=-s_n^2n^2$.
\end{proof}

\begin{theorem}[A lower bound from many zeros; \stU]\label{zfg:thm:ellipselow}
If $p$ is a real polynomial of degree $n\ge1$ with all zeros in $\zfgE_\eps$, $0<\eps\le1$, and $p$ has a Lorentz representation,
then $\zfgdL(p)\ge n/(16\eps)$.
\end{theorem}

\begin{proof}
Write $p=\sum_jb_j(1-X)^j(1+X)^{D-j}$ with $b_j\ge0$, $D=\zfgdL(p)$; $p>0$ on $[-1,1]$, since the open ellipse excludes $\pm1$. If
$\eps\ge1/16$ then $D\ge n$ suffices, so let $v=4\eps<\tfrac14$. For $w\in\C$ let $g(w)\ge0$ be the logarithm of the larger modulus of
the roots of $Z^2-2wZ+1$; $\set{g\le t}$ is the closed ellipse with foci $\pm1$ and semiaxes $\cosh t,\sinh t$. Parametrizing it by
$w=(Z+Z^{-1})/2$, $\abs Z=e^t$, and applying Jensen's formula to $(Z-Z_1)(Z-Z_2)/(2Z)$ gives
\[
 \frac1{2\pi}\int_0^{2\pi}\log\abs{\cos(\theta+it)-w}\dd\theta=-\log2+\max(t,g(w)).
\]
For real $\theta$, $\abs{1\pm\cos(\theta+iv)}=\cosh v\pm\cos\theta$, so $b_j\ge0$ gives $\abs{p(\cos(\theta+iv))}\le(\cosh v)^Dp(\cos\theta/\cosh v)$.
Taking logarithms and integrating, with the identity at $t=v$ on the left and at $t=0$, $w=(\cosh v)\zeta$ on the right,
\[
 \sum_{p(\zeta)=0}\bigl[\max(v,g(\zeta))-g((\cosh v)\zeta)\bigr]\le(D-n)\log\cosh v .
\]
Each term is nonnegative, because $\eps=v/4\le\tanh v$ puts $(\cosh v)\zfgE_\eps$ inside $\set{g\le v}$. A zero $\zeta=x+iy$ with $1-x^2\ge4v^2$
contributes at least $v/2$: with $A=\cosh^2v/\cosh^2(v/2)$ and $B=\eps^2\cosh^2v/\sinh^2(v/2)$ one has $1\le A\le1+2v^2$ and $B\le\tfrac12$, and
$y^2\le\eps^2(1-x^2)$ gives
\[
 \frac{\cosh^2v\,x^2}{\cosh^2(v/2)}+\frac{\cosh^2v\,y^2}{\sinh^2(v/2)}\le A-(A-B)(1-x^2)\le1+2v^2-\tfrac12(1-x^2)\le1,
\]
that is $g((\cosh v)\zeta)\le v/2$. If at least $n/2$ zeros are of this kind, $\log\cosh v\le v^2/2$ gives $D-n\ge n/(2v)$. Otherwise,
after replacing $p(X)$ by $p(-X)$ if needed, at least $n/4$ zeros have $x\ge0$ and $1-x<4v^2$. Differentiating the
representation at $1$, where $b_0>0$, gives $D\ge2p'(1)/p(1)=2\sum_\zeta(1-x)/((1-x)^2+y^2)$, with positive terms; since
$y^2\le2\eps^2(1-x)$, each selected term is at least $1/(1-x+2\eps^2)\ge1/(5v^2)$, so $D\ge n/(10v^2)\ge n/(4v)$. Both cases give
$D\ge n/(4v)=n/(16\eps)$.
\end{proof}

\begin{theorem}[A lower bound from one extreme zero; \stU]\label{zfg:thm:ellipsewidth}
Under the hypotheses of \cref{zfg:thm:ellipselow}, $\zfgdL(p)\ge3/(100\eps^2)$.
\end{theorem}

\begin{proof}
Let $D=\zfgdL(p)$ with $p>0$. For real $s$, $(\cosh s)^Dp(\tanh s)=\sum_ja_je^{(D-2j)s}$ has convex logarithm, whose second derivative is a
variance of $D-2j$ and so nonnegative; differentiating, with $x=\tanh s$,
\[
 D\ \ge\ \sum_{p(\zeta)=0}\Bigl[(1-x^2)\Re\frac1{(x-\zeta)^2}+2x\Re\frac1{x-\zeta}\Bigr]\qquad(-1<x<1),
\]
zeros counted with multiplicity. Reflect so that the leftmost zero $\zeta_0=u+iy$ has $u\le0$, and let $h=1+u$. If $h\le32\eps^2$,
differentiating at $-1$ gives $D\ge-2p'(-1)/p(-1)=2\sum_\zeta\Re(1+\zeta)^{-1}\ge2h/(h^2+y^2)\ge2/(h+2\eps^2)\ge1/(17\eps^2)$, using
$y^2\le\eps^2(1-u^2)\le2\eps^2h$. If $h>32\eps^2$, put $\delta=2\eps\sqrt{1-u^2}$ and $x=u-\delta$. Then $\eps^2\le\delta\le h/2$, so $-1<x<0$ and
$1-x^2\ge(1-u^2)/4$. For any zero $v'+iw'$ with $b=v'-u\ge0$, $w'^2\le\eps^2(1-v'^2)\le\eps^2(1-u^2+2b)\le(\delta+b)^2$, so every term in the
display is nonnegative. For $\zeta_0$, $\abs y\le\delta/2$, and its first term alone gives
$D\ge\tfrac{1-u^2}4\cdot\tfrac{12}{25\delta^2}=\tfrac3{100\eps^2}$.
\end{proof}

\begin{theorem}[Polynomials of small Lorentz degree; \stU]\label{zfg:thm:ellipseup}
For each even $n\ge2$ and $0<\eps\le1$ there is a polynomial $p$ of degree $n$ with rational coefficients and simple zeros, all in
$\zfgE_\eps$, with $\zfgdL(p)\le513/\eps^2$ if $n\eps\le4$ and $\zfgdL(p)\le64n/\eps$ if $n\eps>4$.
\end{theorem}

\begin{proof}
Let $n\eps\le4$, $s=1+\eps^2/4$, $a=\eps/2$ and $P(X)=\cosh(na)+T_n(sX)$. Its zeros $s^{-1}\cos((2j+1)\pi/n+ia)$ lie on an ellipse with
semiaxes $\cosh(a)/s<1$ and $\sinh(a)/s<\eps$, by $\cosh(\eps/2)-1\le\eps^2(\cosh\tfrac12-1)<\eps^2/4$ and $\sinh(\eps/2)\le\eps\sinh\tfrac12<\eps$. On
$[-1,1]$, $P\ge\cosh(na)-1\ge n^2\eps^2/8$ and $\norm{T_n(s\,\cdot)}\le\cosh(n\operatorname{arcosh}(1+\eps^2/4))\le\cosh(n\eps/\sqrt2)\le\cosh(2\sqrt2)<9$. With
$D=\ceil{512/\eps^2}$ and $\eta=2n^2/D\le\tfrac1{16}$, \cref{zfg:lem:bernstein} gives
$\norm{B_D^{-1}F-F}\le36n^2/D<n^2\eps^2/8$ for $F(t)=P(2t-1)$, so $P$ has a positive Lorentz representation of degree
$D\le513/\eps^2$.

Let $n\eps>4$ and $m=2\max(2,\ceil{1/\eps})$, so $2/\eps\le m\le4/\eps<n$. Write $n$ as a sum of even integers $b$ with $m\le b<2m$
($\floor{n/m}$ copies of $m$, the even remainder added to one of them). For each $b$ the proof of \cref{zfg:thm:erdelyi},
which uses only that $b$ is even, shows that $2+T_b((1+b^{-2})X)$ has its zeros in $\zfgE_{2/b}\subseteq\zfgE_\eps$ and a positive
representation of degree $8b^2$. The product has degree $n$ and a positive representation of degree $\sum_b8b^2\le16mn\le64n/\eps$.

In both cases strict positivity of the coefficients in a fixed Lorentz basis and confinement of the zeros to the open
ellipse survive small perturbations of the coefficients, and the discriminant is a nonzero polynomial on coefficient
space, so a rational perturbation has simple zeros.
\end{proof}

\begin{remark}[Exact degrees]\label{zfg:rem:lorentztable}
For $D\ge n$, $P_n$ has a representation of degree $D$ exactly when all coefficients of $(1+z)^{D-n}Q_n(z)$ are nonnegative,
where $Q_n(z)=(n^2)^n(1+z)^nP_n((1-z)/(1+z))$. Integer arithmetic gives $\zfgdL(P_n)=15,56,125,221,495,880,1978,3516$ for
$n=4,8,12,16,24,32,48,64$, with $\zfgdL(P_n)/n^2$ decreasing from $0.9375$ to $0.8584$
(\artifact{erdelyi_ellipse_and_proveit_extensions/01_lorentz_ellipse/exact_degrees.txt}). The product question on the
same page of \cite{Erdelyi1991Estimates} was treated by Ait-Haddou \cite{AitHaddou2015Lorentz}.
\end{remark}

\subsection{Open questions}\label{zfg:sec:open}
\sources{dirichlet_zero_filters; positive-fourier-zeros; kozma_oravecz_half_integer_2026-11-28; ProveIt_New_Results_2026-11-27; erdelyi_ellipse_and_proveit_extensions}

\begin{enumerate}
\item \emph{Progressions and statistics of widths.} For $K=2$, \cref{zfg:thm:threerule} computes $M_2(N,b)$ by a finite
algebraic rule, and closed analytic branches are known for odd-integer ratios $b/N$ with $K=2,3$
(\cref{zfg:thm:three,zfg:thm:four}) and for ratios $d/2$, $d\ge5$ odd, with $K=2$ (\cref{zfg:thm:half}). What are the exact
values for other ratios and for larger $K$, does $K(c-bM_K)$ converge in \cref{rec:thm:wide}, and which polynomials
attain the leading constant? Which wide progressions have rational gaps (\cref{zfg:cor:dichotomy}), and which contact
patterns and algebraic degrees occur for four or more frequencies? At what rate do finite lists of zeros approach
the law of \cref{rec:thm:twolaw}, in particular near the maximum of its phase profile?
\item \emph{Squares.} Does $NM_{N,\infty}$ converge, and to what? \Cref{zfg:thm:longsq} places its limit points in
$[\sqrt3-\tfrac32,\tfrac{21}{62}]$. What is $M_{N,K}$ exactly beyond the short-band regime of \cref{zfg:thm:shortsq}?
\item \emph{Support uncertainty.} What is the least support radius for a given finite set $A$ in
\cref{zfg:thm:support}? Is the linear dependence on the deficit in \cref{zfg:thm:stab} optimal, and with which constant? Is there a sharp
higher-dimensional inequality that takes the directions of the prescribed frequencies into account
(\cref{zfg:rem:dim})?
\item \emph{Lorentz degrees.} What are the best constants in \cref{zfg:thm:lorentz}? Does $\eps^2L_n(\eps)$ have a limiting
profile as $n\to\infty$ and $\eps\downarrow0$ with $n\eps$ tending to a positive constant? Does $\zfgdL(P_n)/n^2$ converge along
$4\mid n$ in \cref{zfg:thm:erdelyi}, and to what?
\end{enumerate}
